\documentclass[11pt]{amsart}

\usepackage{amsmath}
\usepackage{amssymb}
\usepackage{amsthm}
\usepackage{booktabs}
\usepackage{array}
\usepackage{longtable}
\usepackage{hyperref}

\newtheorem{theorem}{Theorem}
\newtheorem{corollary}[theorem]{Corollary}
\newtheorem{proposition}{Proposition}[section]
\newtheorem{lemma}[proposition]{Lemma}
\newtheorem{htheorem}[proposition]{Theorem}
\newtheorem{pcorollary}[proposition]{Corollary}
\theoremstyle{definition}
\newtheorem{definition}[proposition]{Definition}
\theoremstyle{remark}
\newtheorem{remark}[proposition]{Remark}

\numberwithin{equation}{section}
\renewcommand{\theenumi}{\roman{enumi}}
\renewcommand{\labelenumi}{(\theenumi)}
\setlength{\emergencystretch}{2em}
\raggedbottom

\newcommand{\NN}{\mathbb{N}}
\newcommand{\ZZ}{\mathbb{Z}}
\newcommand{\QQ}{\mathbb{Q}}
\newcommand{\RR}{\mathbb{R}}
\newcommand{\CC}{\mathbb{C}}
\newcommand{\DD}{\mathbb{D}}
\newcommand{\TT}{\mathbb{T}}
\newcommand{\fm}{\mathfrak{m}}
\newcommand{\lcm}{\operatorname{lcm}}
\newcommand{\den}{\operatorname{den}}
\newcommand{\Sol}{\operatorname{Sol}}
\newcommand{\Res}{\operatorname{Res}}
\newcommand{\file}[1]{{\def\_{\textunderscore\allowbreak}\texttt{#1}}}
% Visible marker for open items. Every \TODO must be resolved (or the passage deleted) before posting.
\newcommand{\TODO}[1]{\textbf{[TODO: #1]}}

\begin{document}

\hypersetup{pdftitle={Products of two logarithms and arithmetic holonomy bounds},pdfauthor={Shawn Krisman}}
\title[Products of two logarithms]{Products of two logarithms and arithmetic holonomy bounds}

\author{Shawn Krisman}

\subjclass[2020]{Primary 11J72; Secondary 11J82, 34M35, 41A21, 65G20}
\keywords{irrationality, linear independence, products of logarithms, Hermite--Pad\'e approximation, arithmetic holonomy
bounds, monodromy, computer-assisted proof, interval arithmetic}

\begin{abstract}
We prove, assuming the arithmetic holonomy bound of Calegari, Dimitrov and Tang (Theorem~7.0.1 of their preprint
arXiv:2408.15403), that $1$, $\log(1+1/q)$, $\log(1-1/q)$ and $\log(1+1/q)\log(1-1/q)$ are linearly independent over
$\QQ$ for every integer $26\le q\le53$. Together with Hata's theorem for $q\ge54$ this gives the four-term independence
for every $q\ge26$. In particular the seven numbers $\log(1+1/N)\log(1+1/(N+1))$, $25\le N\le31$, are irrational;
irrationality was previously known for $N\ge32$ (Lysov). The proof feeds Lysov's explicit type~II Hermite--Pad\'e forms
into the holonomy bound, and replaces the arithmetic of the forms by an analytic input: the monodromy of a minimal
differential operator of order~13, which we show to be irreducible by interval-arithmetic certificates. Conditionally
on the published form of the Calegari--Dimitrov--Tang argument for products of two logarithms, we also widen their
Theorem~C from $|m/n-1|<10^{-6}$ to $|m/n-1|<3\cdot10^{-6}$, which gives infinitely many new irrational products
$\log(1+1/m)\log(1+1/n)$ with $|m-n|\ge2$, and we show by pooling several pairs that at most one exceptional pair
class exists with $|m/n-1|<5\cdot10^{-5}$, and at most two with $|m/n-1|<10^{-4}$. All certificates are exact or
rigorous ball arithmetic and can be reproduced from the accompanying code at
\url{https://github.com/telaviv/log-products-irrationality}.
\end{abstract}

\maketitle

% ===========================================================================
\section{Introduction}\label{sec:intro}
% ===========================================================================

For an integer $q\ge2$ put
\[
\theta_q=\log\Bigl(1+\frac1q\Bigr)\log\Bigl(1-\frac1q\Bigr).
\]
Hata \cite[Cor.~6.3]{Hat98} showed that $1$, $\log(1+1/q)$, $\log(1-1/q)$, $\theta_q$ are linearly independent over
$\QQ$ for every integer $q\ge54$, with an effective linear independence measure; in particular $\theta_q\notin\QQ$ for
every integer $q\ge54$ \cite[Thm.~0.1]{Hat98}. Lysov \cite{Lys18} found an explicit formula for the common
denominator of the type~II Hermite--Pad\'e approximants to $(\log(1+1/z),\log(1-1/z),\log(1+1/z)\log(1-1/z))$ and
deduced that $\theta_q$ is irrational for every $q\ge33$ \cite{Lys18,Lys19} (\cite{Lys18} is a preprint; we found no
journal version as of 19 November 2026). Calegari, Dimitrov and Tang
\cite[Thm.~C, Thm.~14.0.1]{CDT} proved, by a new arithmetic holonomy bound, that $1$, $\log(1+1/m)$, $\log(1+1/n)$ and
$\log(1+1/m)\log(1+1/n)$ are linearly independent over $\QQ$ whenever $m\ne n$ are integers in $\ZZ\setminus\{-1,0\}$
with $|m/n-1|<10^{-6}$. In \cite[\S15.1.10]{CDT} they suggest forming ``the generating function of the special linear
forms obtained from evaluating a regular sequence of functional Hermite--Pad\'e approximants'' and reading it through
their holonomy bounds. This paper carries out that scheme for Lysov's forms, and pushes the constant $10^{-6}$ of
\cite[Thm.~C]{CDT}.

\subsection*{Dependence on one preprint}
Every result in this paper depends on the preprint \cite{CDT} (arXiv:2408.15403v2, 16 September 2024; still the current
version on 19 November 2026), which to our knowledge has not been refereed. Theorem~\ref{thm:4T} and
Corollary~\ref{cor:F1} use only its general holonomy bound \cite[Thm.~7.0.1]{CDT}, which we apply as stated.
Theorems~\ref{thm:Cprime} and \ref{thm:F2} also use the specific construction of \cite[\S14]{CDT} for products of two
logarithms, together with \cite[Thm.~6.0.2, Cor.~9.0.19, Lemma~12.0.1]{CDT}. That construction contains a step that
is asserted rather than proved in \cite{CDT}, the irreducibility of a certain differential operator
$\mathcal L_{a,b}$, and we inherit it; see Section~\ref{sec:Cprime}. Each theorem below names exactly what it assumes.

\subsection{Results}

\begin{theorem}\label{thm:4T}
Assume \cite[Thm.~7.0.1]{CDT}. For every integer $q$ with $26\le q\le53$, the numbers
\[
1,\qquad\log\Bigl(1+\frac1q\Bigr),\qquad\log\Bigl(1-\frac1q\Bigr),\qquad\log\Bigl(1+\frac1q\Bigr)\log\Bigl(1-\frac1q\Bigr)
\]
are linearly independent over $\QQ$.
\end{theorem}

With Hata's theorem \cite{Hat98} for $q\ge54$, the same holds for every integer $q\ge26$. Precisely,
\cite[Cor.~6.3]{Hat98} states: let $q\ge54$ be an integer; for every $\varepsilon>0$ there is an effectively computable
$H_1=H_1(\varepsilon,q)$ such that
\[
\Bigl|n_0+n_1\log\Bigl(1+\frac1q\Bigr)+n_2\log\Bigl(1-\frac1q\Bigr)+n_3\theta_q+n_4\Lambda\Bigl(\frac1q\Bigr)\Bigr|
\ge H^{-\sigma(q^2)/\tau(q^2)-\varepsilon}
\]
for all $(n_0,\dots,n_4)\in\ZZ^5$ with $H=\max_{1\le m\le4}|n_m|\ge H_1$. Here
$\Lambda(x)=\mathrm{Li}_2((1+x)/2)-\mathrm{Li}_2((1-x)/2)-\log2\cdot\log((1+x)/(1-x))$ \cite[(3.5)]{Hat98}, and
$\sigma,\tau$ are explicit \cite[(6.3), (6.4)]{Hat98}; the bound $q\ge54$ comes from Hata's $\tau(k)>0$ for every integer
$k\ge2909$ (and $53^2<2909\le54^2$), and $\sigma(k)/\tau(k)\to4$ as $k\to\infty$. Taking $n_4=0$ gives the four-term independence; in
fact $1$, $\log(1+1/q)$, $\log(1-1/q)$, $\theta_q$, $\Lambda(1/q)$ are linearly independent over $\QQ$. Hata states
this for integers $q$ only: Corollary~6.3 is the case $k=q^2$ of his Theorem~6.2, the analogous statement at
$1/\sqrt k$ for every integer $k\ge2909$.

Since $\log(1-1/q)=-\log(1+1/(q-1))$, Theorem~\ref{thm:4T} at $q=N+1$ contains the irrationality of the product of two
consecutive logarithms $\log(1+1/N)\log(1+1/(N+1))=-\theta_{N+1}$:

\begin{corollary}\label{cor:F1}
Assume \cite[Thm.~7.0.1]{CDT}. For every integer $N$ with $25\le N\le31$,
\[
\log\Bigl(1+\frac1N\Bigr)\log\Bigl(1+\frac1{N+1}\Bigr)\notin\QQ .
\]
\end{corollary}

These are seven new irrational numbers. By Theorem~\ref{thm:4T} for $32\le N\le52$ and Hata's theorem for $N\ge53$
(and, independently, Lysov's theorem \cite{Lys18} for $N\ge32$), the product is irrational for every $N\ge25$.

The second group of results concerns CDT's two-parameter family. For $m\in\ZZ\setminus\{-1,0\}$ write
$\ell_m=\log(1+1/m)$.

\begin{theorem}\label{thm:Cprime}
Assume the following parts of \cite{CDT} as published: Theorem~6.0.2, Corollary~9.0.19, Lemma~12.0.1 and \S14,
including the irreducibility of the operator $\mathcal L_{a,b}$ and the reduction in the proof of Lemma~14.3.1. Let $m\ne n$ be integers in $\ZZ\setminus\{-1,0\}$ with $|m/n-1|<3\cdot10^{-6}$. Then $1$,
$\ell_m$, $\ell_n$ and $\ell_m\ell_n$ are linearly independent over $\QQ$; in particular $\ell_m\ell_n$ is irrational.
\end{theorem}

\begin{theorem}\label{thm:F2}
Make the assumptions of Theorem~\ref{thm:Cprime}. Call a pair $(m,n)$ of distinct integers in $\ZZ\setminus\{-1,0\}$
\emph{exceptional} if $1,\ell_m,\ell_n,\ell_m\ell_n$ are linearly dependent over $\QQ$, and take classes of pairs
modulo $(m,n)\sim(n,m)\sim(-m-1,-n-1)$. Then:
\begin{enumerate}
\item at most one exceptional class has a representative with $|m/n-1|<5\cdot10^{-5}$;
\item at most two exceptional classes have a representative with $|m/n-1|<10^{-4}$;
\item every exceptional class has $|m/n-1|\ge3\cdot10^{-6}$ in all its representatives (Theorem~\ref{thm:Cprime});
\item in the ranges of (i) and (ii), every exceptional class has $|m-n|\ge2$, by Hata's theorem
\cite[Cor.~6.3]{Hat98}.
\end{enumerate}
\end{theorem}

The relation $(m,n)\sim(-m-1,-n-1)$ comes from $\ell_{-k}=-\ell_{k-1}$: the two pairs give the same four numbers up to
sign. For $1\le N<M$ with $d=M-N$, the most permissive of the four representatives of $\{N,M\}$ is
$(m,n)=(-N-1,-M-1)$, with $|m/n-1|=d/(M+1)$, so
the pair is covered by Theorem~\ref{thm:Cprime} if and only if $d<3\cdot10^{-6}(M+1)$, and by \cite[Thm.~C]{CDT} if and
only if $d<10^{-6}(M+1)$. Consecutive pairs ($d=1$) in either range have $q=N+1>3\cdot10^5$ and are covered by Hata's
theorem. So the new irrational numbers of Theorem~\ref{thm:Cprime} are
\begin{equation}\label{eq:newpairs}
\ell_N\ell_M,\qquad d=M-N\ge2,\qquad 10^{-6}(M+1)\le d<3\cdot10^{-6}(M+1),
\end{equation}
that is, $d\cdot10^6/3<M+1\le d\cdot10^6$: about $\frac23d\cdot10^6$ numbers for each $d\ge2$, infinitely many in all.
For example: for $d=2$, $666{,}664\le N\le1{,}999{,}997$; for $d=3$, $999{,}997\le N\le2{,}999{,}996$; for $d=10$,
$3{,}333{,}323\le N\le9{,}999{,}989$ (exact integer computation over all four representatives,
\file{code/thm3/f7\_newpairs.py}). None of them is a known number in disguise: $(N+1)/N$ is not a perfect power of a
rational number, so $\ell_N$ and $\ell_M$ are multiplicatively independent and $\ell_N\ell_M$ is not of the form
$c\log^2r$. Theorem~\ref{thm:F2} names no new number; it is a statement about the family. For example, all but at most
one of the numbers $\log(1+1/N)\log(1+1/(N+2))$ with $39{,}998\le N\le666{,}663$ are irrational, and all but at most two
with $19{,}998\le N\le666{,}663$.

\subsection{Index conventions}
Four parametrisations are in use in the literature and in our sources. We use the following dictionary throughout.
\begin{center}
\small\setlength{\tabcolsep}{4pt}
\begin{tabular}{lll}
\toprule
convention & consecutive pair & number\\
\midrule
positive form & $(N,N+1)$, $N\ge1$ & $\ell_N\ell_{N+1}$\\
Hata, Lysov, Thm.~\ref{thm:4T} & $q=N+1$ & $\theta_q=-\ell_N\ell_{N+1}$\\
CDT pair $(m,n)$ & $(N,N+1)\sim(-N-2,-N-1)$ & $\ell_m\ell_n$\\
CDT, odd form $(a,b)$ & $(2N+1,2N+3)$ & $a=2m+1$, $b=2n+1$\\
\bottomrule
\end{tabular}
\end{center}
Here $\ell_m=\log\frac{a+1}{a-1}$ when $a=2m+1$.
In Sections~\ref{sec:lysov}--\ref{sec:proof4T} the letter $n$ is the index of Lysov's forms ($P_{j,n}$, $R_{j,n}$);
in Sections~\ref{sec:Cprime}--\ref{sec:F2} $(m,n)$ is a CDT pair, as in the statements above. Lysov's paper writes
$m$ for our $q$, and Zeilberger and Zudilin \cite{ZZ19} write $a=q-1$.

\subsection{What is new}
\begin{itemize}
\item \emph{New irrational numbers.} The seven products $\ell_N\ell_{N+1}$, $25\le N\le31$ (Corollary~\ref{cor:F1});
no method gave their irrationality before. The products \eqref{eq:newpairs} (Theorem~\ref{thm:Cprime}, conditional on
\cite[\S14]{CDT}), together with the four-term independence for each of them.
\item \emph{New linear independence at known irrational numbers.} The four-term independence of
Theorem~\ref{thm:4T} for $26\le q\le32$ is new. For $33\le q\le53$ we know of no explicit statement; ours is the first,
and its proof does not use Lysov's arithmetic. It may follow from Lysov's construction by the classical type~II
criterion, provided his refined denominator estimate covers the numerators $P_{1,n}$, $P_{2,n}$ as well (we have not
checked this). Zeilberger and Zudilin \cite[p.~8]{ZZ19} state that the linear independence of
$1,\log(a/(a+1)),\log((a+1)/(a+2))$, known for $a\ge53$ by Hata, was ``recently improved, though somewhat implicitly,
by Volodya Lysov to $a\ge32$''; that statement concerns the linear forms, not the product.
\item \emph{Nothing new for consecutive pairs in CDT's range or ours}: those are covered by Hata's theorem.
\item \emph{A family statement} (Theorem~\ref{thm:F2}), which names no number.
\item \emph{By-products}: for every $q\ge2$ the characteristic polynomial of Lysov's recurrence is an explicit quartic
(Proposition~\ref{prop:I1}); all local monodromies of the order-34 operator at its finite nonzero singular points are
reflections (Proposition~\ref{prop:exponents}); an explicit form of the factor $\Phi_6$ of \cite[Lemma~14.3.1]{CDT}
and its $2$-adic nonvanishing at odd pairs (Section~\ref{sec:phi6}); a misprinted factor in \cite[(14.2.2), (14.2.3)]{CDT}.
\end{itemize}
A literature search on 19 November 2026 (arXiv, Hexagon, GitHub, Zenodo, general web search) found no other work on
products of two real logarithms of this kind, on the constant $10^{-6}$ of \cite[Thm.~C]{CDT}, or on the
Hata--Lysov family since \cite{Lys18}. The search is not exhaustive; in particular we could not find an abstract of a
2025 talk by Lysov on this subject.

\subsection{Status of the results}\label{sec:status}
The proofs are computer-assisted. Every computation that a proof uses is either an exact identity over $\QQ$ (checked
with SymPy or python-flint) or a rigorous ball-arithmetic enclosure (Arb, through python-flint). The chain has been
checked by AI referees at each decision gate and once as a whole (Section~\ref{sec:provenance}); it has not yet been
read by a human expert, and it is not formalised. In the notation of our working notes (``gate-refereed'': checked by one
adversarial AI referee when the result was accepted; ``chain-refereed'': covered by the final whole-chain AI referee;
``owed'': not yet done), the status labels are:
\begin{itemize}
\item Theorem~\ref{thm:4T} and Corollary~\ref{cor:F1}: \emph{proved, computer-assisted (exact symbolic identities and
Arb certificates), given \cite[Thm.~7.0.1]{CDT} (preprint); gate-refereed once; chain-refereed; human read and Lean
owed.}
\item Theorem~\ref{thm:Cprime}: \emph{proof sketch plus certified computation; conditional on CDT as published
(Thm.~6.0.2, Cor.~9.0.19, Lemma~12.0.1 and \S14, in particular the irreducibility of $\mathcal L_{a,b}$ and the
reduction in the proof of Lemma~14.3.1); $\Phi_6$ factor of Lemma~14.3.1 via Section~\ref{sec:phi6} (proved, exact
computer algebra; not gate-refereed); gate-refereed once; chain-refereed; not in Lean, no human read.}
\item Theorem~\ref{thm:F2}: \emph{proof sketch plus certified computation; conditional on CDT as published (Thm.~6.0.2,
Cor.~9.0.19, Lemma~12.0.1 and \S14, in particular the irreducibility of $\mathcal L_{a,b}$ and the reduction in the
proof of Lemma~14.3.1, reused pair by pair); $\Phi_6$ factor of Lemma~14.3.1 via Section~\ref{sec:phi6} (proved, exact
computer algebra; not gate-refereed); part (iv) adds Hata's theorem \cite[Cor.~6.3]{Hat98} (checked against the paper); Lemma~\ref{lem:D} proved,
given \cite[Lemma~14.2.4]{CDT}; Lemma~\ref{lem:S} ($\varepsilon\le10^{-4}$) and the ray lemma proved;
Lemma~\ref{lem:H} proof sketch (classical); certificates exact (Arb and exact rationals); gate-refereed once;
chain-refereed; not in Lean; no human read.}
\end{itemize}

\subsection{Method}
For Theorem~\ref{thm:4T} we start from Lysov's forms $R_{j,n}(q)=Q_n(q)f_j(q)-P_{j,n}(q)$, $j=1,2,3$, with
$(f_1,f_2,f_3)=(\log(1+1/q),\log(1-1/q),\theta_q)$, and their generating functions $H_j(y)=\sum_nR_{j,n}(q)y^n$. A
rational relation $r_0+a_1f_1+a_2f_2+a_3f_3=0$ makes $-\sum_ja_jH_j$ a power series with rational coefficients of
controlled denominators, which continues analytically to $\CC\setminus[y_2,\infty)$ for an explicit $y_2$. The
holonomy bound \cite[Thm.~7.0.1]{CDT}, applied with the Koebe map onto $\CC\setminus[y_2,\infty)$ and the $14$ functions
$1,H,\vartheta H,\dots,\vartheta^{12}H$ ($\vartheta=y\,d/dy$), gives a contradiction once these $14$ functions are
linearly independent over $\CC(y)$ and $q\ge26$. Their independence follows from the irreducibility of the monodromy of
the minimal operator $L_{13}$ of the $H_j$ (order $13$), which we certify per $q$ from rigorous numerical analytic
continuation, using that the local monodromies at the singular points $y_2,y_3$ are reflections. No arithmetic of the
forms beyond the elementary denominator bound $\lcm(1,\dots,4n)\lcm(1,\dots,2n)$ is needed.

For Theorem~\ref{thm:Cprime} we keep CDT's construction and change only the map: we construct, by optimisation, a
polynomial-exponential map $\psi$ of the disc for which $h\circ\psi$ ($h$ the hauptmodul of $\Gamma_0(2)$) is admissible
for every pair with $|m/n-1|<3\cdot10^{-6}$ at once, and certify the holonomy inequality in its partition form
\cite[Thm.~6.0.2]{CDT} by interval arithmetic. For Theorem~\ref{thm:F2} we put the functions of $K=2$ or $3$
hypothetical exceptional pairs into one holonomy bound, after proving that they are jointly independent.

\subsection*{Outline}
Section~\ref{sec:holonomy} recalls the holonomy bounds. Section~\ref{sec:lysov} treats Lysov's forms: contact orders,
denominators, recurrences and the differential equations. Section~\ref{sec:proof4T} proves Theorem~\ref{thm:4T} and
Corollary~\ref{cor:F1}. Sections~\ref{sec:Cprime} and \ref{sec:F2} prove Theorems~\ref{thm:Cprime} and \ref{thm:F2}.
Section~\ref{sec:computer} lists the computer-assisted steps and how to reproduce them, Section~\ref{sec:limits}
the limits of the method, and Section~\ref{sec:provenance} provenance and the use of AI tools.

% ===========================================================================
\section{The holonomy bounds}\label{sec:holonomy}
% ===========================================================================

We use the notation of \cite[\S6, \S7]{CDT}. For $N\ge1$ write $[1,\dots,N]=\lcm(1,2,\dots,N)$. Let $\fm\ge1$, let
$b=(b_{i,j})_{1\le i\le\fm,\,1\le j\le r}$ be an array of nonnegative reals whose columns have the form
$0=b_{1,j}=\dots=b_{u_j,j}<b_{u_j+1,j}=\dots=b_{\fm,j}$, and let $e\in\NN^\fm$. A formal power series $f_i\in\QQ[[y]]$
has \emph{denominator type} $(b_{i,\cdot},e_i)$ if
\[
f_i(y)=a_{i,0}+\sum_{n\ge1}\frac{a_{i,n}\,y^n}{n^{e_i}\,[1,\dots,b_{i,1}n]\cdots[1,\dots,b_{i,r}n]},\qquad a_{i,n}\in\ZZ
\]
Let $\sigma_i=\sum_jb_{i,j}$ be the row sums (nondecreasing in $i$, by the shape of the columns), and put
\[
\tau(b)=\frac1{\fm^2}\sum_{i=1}^{\fm}(2i-1)\sigma_i ,\qquad\tau(b;e)=\tau(b)+\tau^\sharp(e),
\]
with $\tau^\sharp$ as in \cite[Thm.~6.0.2, Def.~6.0.1]{CDT}; $\tau^\sharp(0)=0$, and $\tau^\sharp(e)$ depends only on
$\fm$, $\sum_ie_i$ and $\max_ie_i$.

For a holomorphic map $\varphi$ from a neighbourhood of the closed unit disc $\overline\DD$ to $\CC$ with $\varphi(0)=0$
put
\[
T(\varphi)=\iint_{\TT^2}\log|\varphi(z)-\varphi(w)|\,d\mu_{\mathrm{Haar}}(z)\,d\mu_{\mathrm{Haar}}(w).
\]

\begin{htheorem}[{\cite[Thm.~7.0.1]{CDT}}]\label{thm:CDT701}
Let $f_1,\dots,f_\fm\in\QQ[[y]]$ be linearly independent over $\QQ(y)$, of denominator types given by $(b,e)$, and let
$\varphi:(\DD,0)\to(\CC,0)$ be holomorphic with $\log|\varphi'(0)|>\max\{\sigma_\fm,\tau(b;e)\}$, such that every
$f_i(\varphi(z))$ is the germ of a meromorphic function on $|z|<1$. Then
\[
\fm\le\frac{T(\varphi)}{\log|\varphi'(0)|-\tau(b;e)} .
\]
If all $f_i$ are known to be holonomic, the hypothesis $\log|\varphi'(0)|>\max\{\sigma_\fm,\tau\}$ may be relaxed to
$\log|\varphi'(0)|>\tau$.
\end{htheorem}

(In \cite{CDT} the map is defined on $\DD$ and $T$ is defined by a limit; we only use maps holomorphic on a
neighbourhood of $\overline\DD$.)

\begin{htheorem}[{\cite[Thm.~6.0.2]{CDT}, partition form}]\label{thm:CDT602}
Let $f_1,\dots,f_\fm$ and $(b,e)$ be as in Theorem~\ref{thm:CDT701}, all $f_i$ holonomic. Let $\varphi_0,\dots,\varphi_l$
be holomorphic maps $(\DD,0)\to(\CC,0)$ with $|\varphi_0'(0)|<\dots<|\varphi_l'(0)|$ and $\log|\varphi_l'(0)|>\tau(b;e)$,
such that every $f_i(\varphi_k(z))$ is meromorphic on $|z|<1$. Let $0=\vartheta_0<\vartheta_1<\dots<\vartheta_l<
\vartheta_{l+1}=\fm$, and let $g:[0,1)\to\RR\cup\{-\infty\}$ be $g(t)=\log|\varphi_k(e^{2\pi i(\fm t-\vartheta_k)/
(\vartheta_{k+1}-\vartheta_k)})|$ for $t\in[\vartheta_k/\fm,\vartheta_{k+1}/\fm)$. Then
\[
\fm\le\frac{\displaystyle\int_0^1\!\!\int_0^1\max(g(s),g(t))\,ds\,dt+\frac1\fm\sum_{k=1}^l\vartheta_k^2\log
\frac{|\varphi_k'(0)|}{|\varphi_{k-1}'(0)|}}{\log|\varphi_l'(0)|-\tau(b;e)} .
\]
\end{htheorem}

With $l=0$ (the \emph{trivial partition}) the numerator is the rearrangement integral
$R(\varphi)=\iint\max(g(s),g(t))\,ds\,dt$ with $g(s)=\log|\varphi(e^{2\pi is})|$. The double integral depends only on
the distribution of $g$, which for a partition is the mixture $\sum_k\frac{\vartheta_{k+1}-\vartheta_k}{\fm}\cdot
(\text{law of }\log|\varphi_k|\text{ on }\TT)$. We use Theorem~\ref{thm:CDT701} in Section~\ref{sec:proof4T} and
Theorem~\ref{thm:CDT602} in Sections~\ref{sec:Cprime} and \ref{sec:F2}.

\begin{lemma}\label{lem:Tunivalent}
If $\varphi$ is univalent on a neighbourhood of $\overline\DD$ and $\varphi(0)=0$, then $T(\varphi)=\log|\varphi'(0)|$.
\end{lemma}

\begin{proof}
The function $(z,w)\mapsto\log|(\varphi(z)-\varphi(w))/(z-w)|$ is pluriharmonic on a neighbourhood of
$\overline\DD^2$ (the quotient is holomorphic and zero-free there, by univalence), so its mean over $\TT^2$ is its value
$\log|\varphi'(0)|$ at $(0,0)$; and $\iint_{\TT^2}\log|z-w|=0$.
\end{proof}

% ===========================================================================
\section{Lysov's forms}\label{sec:lysov}
% ===========================================================================

\subsection{Definitions}
For integers $n\ge0$ and $l$ put
\[
c_{n,l}=\sum_k(-1)^{n+l}\binom{n+l}{l}\binom{n}{l-k}\binom{2n+k}{2k}\binom{n+k}{k}\binom{4n+2k}{2n+k},
\]
the sum over $0\le k\le2n$ with $0\le l-k\le n$ (so $c_{n,l}=0$ unless $0\le l\le3n$), and
\[
Q_n(z)=\sum_{l=0}^{3n}c_{n,l}\,z^{2l}\in\ZZ[z],
\]
an even polynomial of degree $6n$ (Lysov's $Q_{6n}$, \cite[formula~(21) and Lemma~8]{Lys18}; we use the formula as a
definition). Let $f_1(z)=\log(1+1/z)$, $f_2(z)=\log(1-1/z)$ and $f_3=f_1f_2$, as Laurent series at $z=\infty$. Let
$P_{j,n}$ be the polynomial part of $Q_nf_j$, and $R_{j,n}=Q_nf_j-P_{j,n}=O(z^{-1})$. Since $Q_n$ is even and
$f_2(z)=f_1(-z)$, we have $P_{2,n}(z)=P_{1,n}(-z)$. Explicitly, with $A_N=\sum_{i=1}^N(-1)^{i+1}/i$,
\begin{equation}\label{eq:P}
\begin{aligned}
P_{1,n}(z)&=\sum_lc_{n,l}\sum_{m=1}^{2l}\frac{(-1)^{m+1}}{m}z^{2l-m},\\
P_{3,n}(z)&=\sum_lc_{n,l}\sum_{m=1}^{l}\Bigl(-\frac{A_{2m-1}}m\Bigr)z^{2(l-m)},
\end{aligned}
\end{equation}
because $[z^{-2m}]f_3=-A_{2m-1}/m$ and the odd negative powers of $z$ do not occur in $f_3$.

The functions are Cauchy transforms: for $|z|>1$,
\[
f_1(z)=\int_{-1}^0\frac{dx}{z-x},\qquad f_2(z)=-\int_0^1\frac{dx}{z-x},\qquad
f_3(z)=\int_{-1}^1\frac{\omega(x)\,dx}{z-x},
\]
with $\omega$ odd and $\omega(x)=-\log(1+1/x)$ on $(0,1)$. Write $\mu_j$ for the three (signed) measures; their total
variations are $1,1,4\log2$. Then $R_{j,n}(z)=\int Q_n(x)\,d\mu_j(x)/(z-x)$.

\subsection{Contact orders}

\begin{lemma}[Lysov]\label{lem:contact}
For all $n\ge0$ and $j=1,2,3$: $\int x^kQ_n(x)\,d\mu_j(x)=0$ for $0\le k<2n$. Equivalently $R_{j,n}(z)=O(z^{-2n-1})$ as
$z\to\infty$.
\end{lemma}

This is \cite[Lemmas~7 and 8]{Lys18}. Since \cite{Lys18} is an unrefereed preprint we give the argument in full; it
follows Lysov's.

\begin{proof}
Put $\mathfrak d=\frac1x\frac d{dx}$ and $\mathfrak D=\frac d{dx}\circ\frac1x$, so that $\mathfrak dx^m=mx^{m-2}$,
$\mathfrak Dx^m=(m-1)x^{m-2}$, both operators commute with $x\mapsto-x$, and for smooth $f,g$ on $(0,1]$
\[
\int_\epsilon^1f\,\mathfrak dg=\Bigl[\frac{fg}x\Bigr]_\epsilon^1-\int_\epsilon^1(\mathfrak Df)g,\qquad
\int_\epsilon^1f\,\mathfrak Dg=\Bigl[\frac{fg}x\Bigr]_\epsilon^1-\int_\epsilon^1(\mathfrak df)g .
\]

\emph{Rodrigues form.} We have
\[
\begin{gathered}
Q_n=\kappa_n\,\mathfrak d^n\bigl[(x^2-1)^nx^{2n}\,\tilde U\bigr],\qquad\kappa_n=\frac1{(2n)!\,(n!)^2},\\
\tilde U=\mathfrak d^n[x^{2n}U],\qquad U=\mathfrak D^{2n}[(x^2-1)^{2n}x^{4n}].
\end{gathered}
\]
In the variable $w=x^2$ ($\mathfrak d=2\,d/dw$, $\mathfrak D^{2n}=x\,\mathfrak d^{2n}x^{-1}$, $D=d/dw$) this is the pair of
identities
\[
\begin{aligned}
&\text{(i)}&Q_n(w)&=\frac1{n!}D^n\bigl[w^n(w-1)^nV_n(w)\bigr],\\
&\text{(ii)}&V_n(w)&=\frac{16^n}{(2n)!\,n!}D^n\bigl[w^{n+1/2}D^{2n}[w^{2n-1/2}(w-1)^{2n}]\bigr],
\end{aligned}
\]
with $V_n(w)=\sum_k(-1)^kB_kw^k$, $B_k=\binom{2n+k}{2k}\binom{n+k}{k}\binom{4n+2k}{2n+k}$. For (i),
$\frac1{n!}D^n[w^{n+k}(w-1)^n]=\sum_j\binom nj(-1)^{n-j}\binom{n+k+j}{n}w^{k+j}$; multiplying by $(-1)^kB_k$ and
collecting $w^l$ gives $\sum_{k+j=l}(-1)^{n+l}\binom{n+l}{l}\binom nj B_k=c_{n,l}$. For (ii), the coefficient of $w^k$ on
the right is $\frac{16^n}{(2n)!\,n!}\binom{2n}k(-1)^k\frac{\Gamma(2n+k+\frac12)}{\Gamma(k+\frac12)}\frac{(n+k)!}{k!}$,
which equals $(-1)^kB_k$ by the Legendre duplication formula for $0\le k\le2n$ (the ratio simplifies to $1$; checked
symbolically in \file{i1\_contact.py}), and both sides vanish for $k>2n$. The factors $U$, $x^{2n}U$, $\tilde U$ are even
polynomials, and $x^{2n}U$ and $(x^2-1)^nx^{2n}\tilde U$ vanish to order $\ge2n$ at $0$.

\emph{Moments.} By the Cauchy-transform representations, $[z^{-s-1}]R_{j,n}=\int x^sQ_n\,d\mu_j$. (For $f_3$ this is the
identity $[z^{-2m}]f_3=-A_{2m-1}/m=-2\int_0^1x^{2m-1}\log(1+1/x)\,dx$, which follows from
$\int_0^1x^{N-1}\log(1+x)\,dx=A_N/N$ for even $N$ and $\int_0^1x^{N-1}(-\log x)\,dx=1/N^2$; odd negative powers vanish on
both sides.) Since $Q_n$ is even, the lemma is equivalent to
\[
\begin{aligned}
&\text{(O1)}&\int_0^1x^sQ_n\,dx&=0\quad(0\le s\le2n-1),\\
&\text{(O2)}&\int_0^1x^sQ_n\,\ell\,dx&=0\quad(s\text{ odd},\ s\le2n-1),
\end{aligned}
\]
where $\ell(x)=\log(1+1/x)$. Let $h=x^s$ or $h=x^s\ell$ accordingly.

\emph{Step 1.} Integrating by parts $n$ times against $\mathfrak d$ gives
$\int_0^1Q_nh=\pm\kappa_n\int_0^1(x^2-1)^nx^{2n}\tilde U\cdot r\,dx$ with $r=\mathfrak D^nh$. The boundary terms vanish:
$\mathfrak d^{n-1-j}[(x^2-1)^nx^{2n}\tilde U]$ is divisible by $(x^2-1)^{j+1}$ and vanishes to order $2j+2$ at $0$, while
$\mathfrak D^jh=O(x^{-2j}|\log x|)$ at $0$.

\emph{Step 2.} If $h=x^s$ with $s$ odd, then $r=0$ ($\mathfrak D^n$ kills odd polynomials of degree $\le2n-1$). If $h=x^s$
with $s$ even, $r=c\,x^{s-2n}$. If $h=x^s\ell$, then by induction $\mathfrak D^j[x^s\ell]=\mathfrak D^j[x^s]\,\ell+
(\text{rational})$, since $\ell'=-1/(x(x+1))$; as $\mathfrak D^nx^s=0$, $r$ is rational with poles only at $0$ (order
$\le2n$) and $-1$ (order $\le n$), and $r=O(x^{-2})$ at $\infty$. Its odd part is $\mathfrak D^n[x^s\ell_{\mathrm{even}}]$
with $\ell_{\mathrm{even}}=\frac12\log(1-x^{-2})=-\sum_{j\ge1}x^{-2j}/(2j)$; here $\mathfrak D^nx^{s-2j}=0$ for
$1\le j\le(s-1)/2$, and the remaining powers are $\le x^{-2n-1}$, so the odd part of $r$ is $O(x^{-2n-1})$. Put
$B=x^{2n}(x^2-1)^nr$. In both nontrivial cases $B$ is a polynomial divisible by $(x-1)^n$, of degree $\le4n-2$, whose odd
part $x^{2n}(x^2-1)^nr_{\mathrm{odd}}$ has degree $\le2n-1$; and $\int_0^1Q_nh=\pm\kappa_n\int_0^1\tilde U B\,dx$.

\emph{Step 3.} Integrating by parts $n$ times against $\mathfrak d$ ($\tilde U=\mathfrak d^n[x^{2n}U]$) gives
$\int_0^1\tilde UB=\pm\int_0^1x^{2n}U\,\mathfrak D^nB$; the boundary terms vanish because $\mathfrak D^jB$ is divisible by
$(x-1)^{n-j}$ and is $O(x^{-2j})$, while $\mathfrak d^{n-1-j}[x^{2n}U]$ vanishes to order $2j+2$ at $0$.

\emph{Step 4.} Integrating by parts $2n$ times against $\mathfrak D$ ($U=\mathfrak D^{2n}E$, $E=(x^2-1)^{2n}x^{4n}$) gives
$\int_0^1x^{2n}U\,\mathfrak D^nB=\pm\int_0^1\mathfrak d^{2n}[x^{2n}\mathfrak D^nB]\cdot E$; the boundary terms vanish because
$\mathfrak D^{2n-1-j}E$ is divisible by $(x^2-1)^{j+1}$ and vanishes to order $2j+2$ at $0$, while $\mathfrak d^j$ of a
polynomial is $O(x^{-2j})$.

\emph{Step 5.} $x^{2n}\mathfrak D^nB$ is a polynomial of degree $\le4n-2$ whose odd part $x^{2n}\mathfrak D^n[B_{\mathrm
{odd}}]$ vanishes ($\deg B_{\mathrm{odd}}\le2n-1$). An even polynomial of degree $\le4n-2$ is killed by $\mathfrak d^{2n}$,
so every integral is $0$.

As sanity checks, \file{i1\_contact.py} verifies (i) and (ii) exactly for small $n$, the moments of $f_3$ for $m\le30$, and
the contact orders exactly for $n\le24$.
\end{proof}

\begin{pcorollary}\label{cor:size}
$|Q_n(x)|\le(2n+1)2^{20n}$ on $[-1,1]$, and for every integer $q\ge2$ and $j=1,2,3$,
\[
|R_{j,n}(q)|\le(2n+1)\Bigl(\frac{2^{20}}{q^2}\Bigr)^n\frac{|\mu_j|}{q-1}.
\]
\end{pcorollary}

\begin{proof}
We have
\[
\begin{aligned}
\sum_l|c_{n,l}|&\le\sum_k\max_l\binom{n+l}l\,2^n\binom{2n+k}{2k}\binom{n+k}k\binom{4n+2k}{2n+k}\\
&\le(2n+1)2^{4n+n+4n+3n+8n}.
\end{aligned}
\]
By Lemma~\ref{lem:contact}, $R_{j,n}(q)=q^{-2n}\int Q_n(x)x^{2n}\,d\mu_j(x)/(q-x)$.
\end{proof}

\subsection{Denominators}
The following lemma is \cite[Lemma~9]{Lys18} at prime-power level; we give the proof in full, so that \cite{Lys18} is
used only for the definition of $c_{n,l}$.

\begin{lemma}\label{lem:den}
For every $n\ge1$ and $j=1,2,3$, the coefficients of $P_{j,n}$ have denominators dividing $[1,\dots,4n]\,[1,\dots,2n]$.
\end{lemma}

\begin{proof}
Write $D_k=[1,\dots,k]$ and fix a prime $p$. We use Kummer's theorem in the form
$v_p\binom{A}{B}=\sum_{i\ge1}(\lfloor A/p^i\rfloor-\lfloor B/p^i\rfloor-\lfloor(A-B)/p^i\rfloor)$, each term being $0$
or $1$. In every nonzero summand of $c_{n,l}$ we have $l-n\le k\le2n$.

\emph{$P_1$ and $P_2$.} By \eqref{eq:P} the denominators of $P_{1,n}$ divide the numbers $m\le2l\le6n$, with
coefficient $c_{n,l}$. If $p\le2n$, then $v_p(m)\le\lfloor\log_p6n\rfloor\le\lfloor\log_p4n\rfloor+1\le
v_p(D_{4n}D_{2n})$. If $2n<p\le4n$, then $p^2>6n$, so $v_p(m)\le1=v_p(D_{4n})$. If $p>4n$ and $p\mid m$, then
$4n<p<2l$ ($p$ is odd). Then for every summand of $c_{n,l}$: $2n+k\le4n<p$ and $p<2l\le2n+2k\le4n+2k\le8n<2p$, so
$\lfloor(4n+2k)/p\rfloor-2\lfloor(2n+k)/p\rfloor=1$ and $p\mid\binom{4n+2k}{2n+k}$. Hence $p\mid c_{n,l}$, and
$c_{n,l}/m$ is $p$-integral. (The signs $(-1)^{m+1}$ play no role.) $P_{2,n}(z)=P_{1,n}(-z)$.

\emph{$P_3$.} By \eqref{eq:P} the term $c_{n,l}A_{2m-1}/m$ ($m\le l\le3n$) has $v_p\ge v_p(c_{n,l})-e_1-e_2$ with
$e_1=v_p(m)$ and $e_2=\lfloor\log_p(2m-1)\rfloor$. Now $e_1\le\lfloor\log_p2n\rfloor+\delta_1$, where $\delta_1=1$ if
some $p^{a}\mid m$ with $p^a\in(2n,3n]$ and $\delta_1=0$ otherwise (a higher excess is impossible since
$2p^a>3n$); and $e_2\le\lfloor\log_p4n\rfloor+\delta_2$, where $\delta_2=1$ if some $p^b\in(4n,2m-1]$ and $\delta_2=0$
otherwise (again $2p^b>6n$). If $\delta_1+\delta_2>0$ then $m>2n$, so $l>2n$ and $k\ge l-n>n$ in every summand. If
$\delta_1=1$: $n<p^a$, $k\le2n<p^a$ and $n+k\ge l\ge m\ge p^a$, so $p\mid\binom{n+k}{k}$. If $\delta_2=1$:
$2n+k\le4n<p^b<2m\le2l\le2n+2k\le4n+2k$, so $p\mid\binom{4n+2k}{2n+k}$. These are different binomial factors, so
$v_p(c_{n,l})\ge\delta_1+\delta_2$, and the term has $v_p\ge-v_p(D_{2n}D_{4n})$.
\end{proof}

(An exact check of the divisibility for $n\le120$, and independently for selected $n\le81$, agrees.)

\subsection{The recurrence}

\begin{proposition}\label{prop:I1}
There are polynomials $p_0,\dots,p_4\in\ZZ[n,w]$, of degree $34$ in $n$ and of degrees $3,4,5,6,3$ in $w$, such that
for every integer $n\ge0$ and every $z\in\CC$, with $w=z^2$,
\[
\sum_{i=0}^4p_i(n,z^2)\,X_{n+i}(z)=0\qquad\text{for }X\in\{Q,P_1,P_2,P_3\},
\]
hence also for $X=R_j$. Moreover
\[
\sum_{j=0}^4\mathrm{lc}_n(p_j)(w)\,\lambda^j=2^{59}(8w+1)(9w-1)(25w-16)\cdot\chi_w(-\lambda),
\]
where
$\chi_w(\lambda)=\lambda^4+(16384w^3-20480w^2+5248w-64)\lambda^3+(135168w^2+159744w+1536)\lambda^2+(294912w-16384)\lambda+
65536$.
\end{proposition}

\begin{proof}[Proof (computer-assisted, exact)]
\emph{Two relations for $c_{n,l}$.} There are explicit polynomials $a_0,a_1,a_2$ and $\alpha,\beta_{a,i}$ with
\[
\sum_{m=0}^2a_m(n,l)c_{n,l+m}=0,\qquad
\alpha(n)c_{n+1,l}=\sum_{a=0}^3\sum_{i=0}^3\beta_{a,i}(n)(l-i)^ac_{n,l-i},
\]
$a_0=-(l-3n)(l+n+1)(l+n+2)(2l+2n+1)$, $a_2=-(l+2)^3(2l+3)$, $\alpha=-(n+1)^3(2n+1)^2(4n+1)(4n+3)$ (the others are in
\file{runs/i1\_relations.json}). They are proved by creative telescoping in $k$: with explicit certificates $G_1,G_2$
(hypergeometric terms given by factorial formulas), the summand identities
$\sum_ma_mF(n,l+m,k)=G_1(n,l,k+1)-G_1(n,l,k)$ and its analogue hold at every integer point of the relevant regions.
Over a common factorial base each reduces to one polynomial identity in $\QQ[n,l,k]$, which is identically zero; at
$k=-1$ the certificates vanish (factor $k^2$), they vanish for $k>2n+2$, so the sums over $k$ telescope; the finitely
many boundary cases are checked directly (\file{i1\_certs.py}).

\emph{Order-4 operator.} In generating-function form the relations say $D_nQ_n=0$, with $D_n$ an operator of order $4$
in $w\,d/dw$ whose leading coefficient $-2(w-1)^2$ does not depend on $n$, and $\alpha(n)Q_{n+1}=E_nQ_n$ with $E_n$ of
order $3$. Reducing $E_{n+i}\cdots E_n$ modulo $D_n$ on the right expresses $Q_{n+i}$, $i=0,\dots,4$, through
$Q_n,\dots,(w\,d/dw)^3Q_n$; the cofactor vector of the resulting $4\times5$ matrix, made primitive, is
$(p_0,\dots,p_4)$, and $\sum_ip_iQ_{n+i}=0$ is re-verified exactly (\file{i1\_opalg.py}).

\emph{The numerators.} $\sum_ip_i(n,z^2)P_{j,n+i}=-\sum_ip_i(n,z^2)R_{j,n+i}$. The left side is a polynomial, and by
Lemma~\ref{lem:contact} the right side is $O(z^{5-2n})$; so both vanish for $n\ge3$, and $n=0,1,2$ are checked directly.

The characteristic polynomial identity is exact (\file{i1\_extras.py}).
\end{proof}

For $q\ge2$ the prefactor $2^{59}(8w+1)(9w-1)(25w-16)$ does not vanish at $w=q^2$, so the characteristic roots of the
recurrence are the roots $\lambda$ of $\chi_{q^2}(-\lambda)$. For $2\le q\le60$ they are real, positive and distinct, and
$1/\lambda_1<2^{-20}q^2$ (Arb root isolation, \file{fourterm/check\_y1.py}); we order them
$\lambda_1>\lambda_2>\lambda_3>\lambda_4>0$ and put $y_k=1/\lambda_k$, so that $0<y_1<y_2<y_3<y_4$. For example, for
$q=26$: $y_1\approx1.98\cdot10^{-13}$, $y_2\approx156.37$, $y_3\approx182.38$, $y_4\approx2703$. (For every $q\ge2$ the
roots are distinct, since the discriminant of $c_{34}$ below does not vanish at $w=q^2$; we use their reality and order
only for $26\le q\le53$.)

\subsection{The differential equations}
Fix an integer $q$ with $2\le q\le60$ and put
\[
H_j(y)=\sum_{n\ge0}R_{j,n}(q)\,y^n\quad(j=1,2,3),\qquad H_a=\sum_{j=1}^3a_jH_j\quad(a\in\CC^3).
\]
By Corollary~\ref{cor:size}, each $H_j$ is holomorphic on $|y|<2^{-20}q^2$, a disc containing $0$ and $y_1$. Let
$\vartheta=y\,d/dy$ and
\[
L=\sum_{j=0}^4y^{4-j}p_j(\vartheta-j,q^2),
\]
a differential operator of order $34$. The coefficient of $y^{N+4}$ in $L\sum_nX_ny^n$ is $\sum_jp_j(N)X_{N+j}$, so by
Proposition~\ref{prop:I1}, $LH_j=E_j$ with $E_j$ a polynomial of degree $\le3$. The $\vartheta^{34}$-coefficient of $L$
is $c_{34}(y)=\sum_jy^{4-j}\mathrm{lc}_n(p_j)$, which is a constant times $(8w+1)(9w-1)(25w-16)\,y^4\chi_w(-1/y)$; its
roots are $y_1,\dots,y_4$, and $c_{34}(0)\ne0$. Hence the singular points of $L$ are $0,y_1,\dots,y_4,\infty$, and $L$ has
no apparent singularities.

\begin{proposition}\label{prop:continuation}
For every $a\in\CC^3$, $H_a$ continues analytically to a single-valued holomorphic function on $\CC\setminus[y_2,\infty)$.
\end{proposition}

\begin{proof}
$H_a$ is holomorphic on the disc $|y|<2^{-20}q^2$, which contains $0$ and $y_1$, and satisfies $LH_a=E_a$; solutions of a
linear differential equation continue along every path avoiding its singular points. The fundamental group of
$\CC\setminus([y_2,\infty)\cup\{0,y_1\})$ is generated by loops around $0$ and $y_1$ that lie in that disc, where $H_a$ is
single-valued; so the continuation is single-valued on $\CC\setminus([y_2,\infty)\cup\{0,y_1\})$ and, being holomorphic
near $0$ and $y_1$, extends holomorphically to $\CC\setminus[y_2,\infty)$.
\end{proof}

\subsection{The minimal operator}
For each $q$ with $26\le q\le53$ we computed an operator of recurrence order $25$ and $n$-degree $13$, common to the four
sequences $Q,P_1,P_2,P_3$ at $z=q$ (modular nullspaces, Chinese remaindering and rational reconstruction; checked
exactly for $n\le60$). Its associated differential operator $L_{13}=\sum_{j=0}^{25}y^{25-j}\tilde p_j(\vartheta-j)$ has
order $13$. Its leading coefficient has degree $25$: its roots are $y_1,\dots,y_4$ and $21$ apparent singular points.

\begin{proposition}\label{prop:I1prime}
For each $q$ with $26\le q\le53$ the order-$25$ recurrence holds for all $n\ge0$ on $Q,P_1,P_2,P_3$ at $z=q$. Hence
$L_{13}H_a=E'_a$ with $E'_a$ a polynomial, for every $a\in\CC^3$.
\end{proposition}

\begin{proof}[Proof (computer-assisted, exact)]
Fraction-free right division in $\ZZ[n][S]$ ($S$ the shift) gives $\mu(n)\cdot\mathrm{Rec}_{25}=\tilde K\cdot L_4$ with
zero remainder, where $L_4=\sum_ip_i(n,q^2)S^i$ and $\mu(n)$ is a nonzero integer times
$\prod_{t=0}^{21}p_4(n+t,q^2)$. The polynomial $p_4(\cdot,q^2)$ has no root at a nonnegative integer, so $\mu(n)\ne0$
for $n\ge0$, and Proposition~\ref{prop:I1} gives the claim (\file{i1\_rec25.py}).
\end{proof}

% ===========================================================================
\section{Proof of Theorem~\ref{thm:4T}}\label{sec:proof4T}
% ===========================================================================

Fix $q$ with $26\le q\le53$. Write $f=(f_1(q),f_2(q),f_3(q))$, $P_n=(P_{1,n}(q),P_{2,n}(q),P_{3,n}(q))$ and
$v_n=(Q_n(q),P_n)\in\QQ^4$.

\subsection{Reduction}
Suppose $r_0+a_1f_1+a_2f_2+a_3f_3=0$ with $(r_0,a)\in\ZZ^4\setminus\{0\}$. Then $a\ne0$, and
\[
G(y):=\sum_{n\ge0}\bigl(r_0Q_n(q)+a\cdot P_n\bigr)y^n
=\sum_n\bigl(Q_n(q)(r_0+a\cdot f)-a\cdot R_n\bigr)y^n=-H_a(y),
\]
where $R_n=(R_{1,n}(q),R_{2,n}(q),R_{3,n}(q))$.
By Lemma~\ref{lem:den} (and $Q_n(q)\in\ZZ$), $G\in\QQ[[y]]$ has constant term $r_0\in\ZZ$ and denominator type
$[1,\dots,4n][1,\dots,2n]$; so do $\vartheta^jG$ for all $j$ (the coefficients are multiplied by $n^j$). By
Proposition~\ref{prop:continuation}, $G$ continues to $\CC\setminus[y_2,\infty)$.

Let $\varphi(w)=4y_2w/(1+w)^2$, the Koebe map of $\DD$ onto $\CC\setminus[y_2,\infty)$, and $\varphi_r(w)=\varphi(rw)$
for $0<r<1$. Consider the $\fm=14$ functions $1,G,\vartheta G,\dots,\vartheta^{12}G$, with type rows $(0,0)$ for the
constant $1$ and $(4,2)$ for the others, $e=0$. Then $\sigma_i=6$ for $i\ge2$ and
$\tau=\frac1{196}\sum_{i=2}^{14}(2i-1)\cdot6=6\cdot\frac{195}{196}$. All pullbacks by $\varphi_r$ are holomorphic on
a neighbourhood of $\overline\DD$, and $\varphi_r$ is univalent, so $T(\varphi_r)=\log|\varphi_r'(0)|=L_q+\log r$ by
Lemma~\ref{lem:Tunivalent}, where
\[
L_q=\log(4y_2).
\]
If the $14$ functions are linearly independent over $\QQ(y)$, Theorem~\ref{thm:CDT701} gives
$14\le(L_q+\log r)/(L_q+\log r-\tau)$ for every $r<1$ close to $1$, hence $14\le L_q/(L_q-\tau)$, that is
\[
M_{14}(q):=\tfrac{13}{14}L_q-6\cdot\tfrac{195}{196}\le0.
\]
(The condition $L_q>\max\{6,\tau\}=6$ of Theorem~\ref{thm:CDT701} holds.) But $M_{14}(q)>0$ for every
$q$ with $26\le q\le53$; the Arb values are in Table~\ref{tab:4T}, for example $M_{14}(26)=0.0092251638\pm5\cdot10^{-12}$ and
$M_{14}(53)=1.3688032\pm3\cdot10^{-8}$ (\file{budget\_4t.py}; $y_2$ is a root of the exact characteristic polynomial of
Proposition~\ref{prop:I1}). This contradiction proves Theorem~\ref{thm:4T}, given:

\begin{proposition}\label{prop:indep}
Let $a\in\CC^3\setminus\{0\}$. The $14$ functions $1,H_a,\vartheta H_a,\dots,\vartheta^{12}H_a$ are linearly
independent over $\CC(y)$, hence over $\QQ(y)$.
\end{proposition}

The rest of this section proves Proposition~\ref{prop:indep}. For $a=(0,0,1)$ it is the case used for
Corollary~\ref{cor:F1}, which we also certified separately (Remark~\ref{rem:F1cert}).

\subsection{An independence lemma}
Fix a base point $y_*\in(y_1,y_2)$ which is a regular point of $L_{13}$, and consider germs at $y_*$; $\gamma$ runs
over closed paths at $y_*$ avoiding all singular points of $L_{13}$. Let
\[
U_a=\operatorname{span}_\CC\{(\gamma-1)H_a\},\qquad S=\Sol(L_{13})\ (\dim S=13).
\]
By Proposition~\ref{prop:I1prime}, $U_a\subseteq S$. $U_a$ is invariant under monodromy, since
$\delta(\gamma-1)H_a=(\delta\gamma-1)H_a-(\delta-1)H_a$.

\begin{lemma}\label{lem:indep}
If $\dim U_a=u$, then $1,H_a,\vartheta H_a,\dots,\vartheta^{u-1}H_a$ are linearly independent over $\CC(y)$.
\end{lemma}

\begin{proof}
Suppose $c_0+KH_a=0$ with $K=\sum_{j<u}c_{j+1}\vartheta^j$, $c_i\in\CC(y)$. Continuation along $\gamma$ (perturbed off
the poles of the $c_i$) gives $c_0+K(\gamma H_a)=0$, so $K$ kills $(\gamma-1)H_a$ for every $\gamma$, hence the
$u$-dimensional space $U_a$. A nonzero operator of order $\le u-1$ cannot do that, so $K=0$, and then $c_0=0$.
\end{proof}

So Proposition~\ref{prop:indep} follows from $U_a=S$, which in turn follows from $U_a\ne0$
(Lemma~\ref{lem:nonzero}) and the irreducibility of $S$ under the monodromy (Proposition~\ref{prop:irred}).

\subsection{Local exponents}

\begin{proposition}\label{prop:exponents}
For every integer $q\ge2$, each of the four roots $y_k$ of $c_{34}$ is a regular singular point of $L$ with exponents
$0,1,\dots,32$ and $\tfrac12$. Every local solution of $Lf=E$ ($E$ holomorphic at $y_k$) has the form
$h_0(t)+c\,t^{1/2}h(t)$, $t=y-y_k$, with $h_0,h$ holomorphic and $h(0)\ne0$. In particular the local monodromy of the
affine solution space of $Lf=E$ at $y_k$ is a reflection, and a solution is single-valued around $y_k$ if and only if it
is holomorphic there.
\end{proposition}

\begin{proof}[Proof (computer-assisted, exact)]
Let $c_{34},c_{33}$ be the $\vartheta^{34}$- and $\vartheta^{33}$-coefficients of $L$. In $d/dy$ form,
$a_{34}=y^{34}c_{34}$ has a simple zero at $y_k$ (the discriminant of $c_{34}$ is a constant times
$w^3(w-1)^2(8w+1)^8(9w-1)^8(25w-16)^8\ne0$) and $a_{33}(y_k)=y_k^{33}c_{33}(y_k)$. So $y_k$ is regular singular with
exponents $0,\dots,32$ and $\beta_k=33-c_{33}(y_k)/(y_kc_{34}'(y_k))$. The pseudo-remainder of $2c_{33}-65\,y\,c_{34}'$
by $c_{34}$ in $\ZZ[w][y]$ is identically zero, so $\beta_k=\frac12$ for all $k$ and all $q$ (cross-check:
$\Res_y(c_{34},(33-b)yc_{34}'-c_{33})$ is a constant times $w^3(w-1)^2(8w+1)^{10}(9w-1)^{10}(25w-16)^{10}(2b-1)^4$). In
the power-series recursion at $y_k$ the pivot vanishes only at exponent $\frac12$, never at an integer, which gives $33$
holomorphic homogeneous solutions, a holomorphic particular solution and one solution $t^{1/2}h(t)$
(\file{fourterm/exponents.py}).
\end{proof}

These statements concern $L$, not $L_{13}$ (which has its own exponents and apparent points); they are transferred to
$S$ only after $S\subseteq\Sol(L)$ has been certified (step (a) below).

\subsection{No combination is single-valued}

\begin{lemma}\label{lem:nonzero}
For every $a\in\CC^3\setminus\{0\}$, $U_a\ne0$.
\end{lemma}

\begin{proof}
Suppose $U_a=0$, so that $H_a$ extends to a single-valued holomorphic function on $\CC\setminus\{0,y_1,\dots,y_4\}$.
\begin{enumerate}
\item It is holomorphic at $0$ and $y_1$ (where it equals its power series) and at $y_2,y_3,y_4$ by
Proposition~\ref{prop:exponents}. So $H_a$ is entire.
\item $E_a$ is a polynomial of degree $\le3$, so $H_a$ is killed by $L'=\vartheta(\vartheta-1)(\vartheta-2)
(\vartheta-3)L$. With $u=1/y$, the $\vartheta_u$-leading coefficient of $L'$ at $u=0$ is $\pm\mathrm{lc}(p_0)\ne0$, so
$\infty$ is a regular singular point. A single-valued solution of a Fuchsian equation near $\infty$ has moderate growth,
hence is meromorphic at $\infty$: $H_a$ is a polynomial.
\item $Ly^D=\sum_jp_j(D-j)y^{D+4-j}$, so if $H_a\not\equiv0$ has degree $D\ge0$, the coefficient of $y^{D+4}$ in $LH_a$
is $p_0(D,q^2)$ times the leading coefficient of $H_a$; $p_0(\cdot,q^2)$ has no nonnegative integer root for
$2\le q\le60$. Since $\deg E_a\le3<D+4$, $H_a\equiv0$, that is, $(a\cdot f)Q_n-a\cdot P_n=0$ for all $n$.
\item The Casoratian $\det[v_0,v_1,v_2,v_3]$ equals, as a function of $z=q$,
\[
\begin{aligned}
-\frac{z}{291060}\bigl(&7427354003333880z^6-4691165040913890z^4\\&-141637476029121z^2+63093459176677\bigr),
\end{aligned}
\]
which has no integer root other than $0$. So $v_0,\dots,v_3$ are independent and $(a\cdot f,-a)=0$, i.e.\ $a=0$.
\qedhere
\end{enumerate}
\end{proof}

\subsection{Irreducibility of the monodromy}\label{sec:irred}

\emph{The data.} For each $q$, the script \file{mono13.py} computes rigorous enclosures of the monodromy of $g=y^5H_3$
(multiplication by $y^5$ is equivariant, and shifts the exponents of $L_{13}$ at $0$ to make the Taylor recursion
well-conditioned) on the $14$-dimensional affine module spanned by $g$ and $y^5S$. Coordinates: the $13$ scaled Taylor
coefficients at the base point $y_*=\operatorname{round}(0.42y_2)$ and an affine coordinate. The generators are: $A$,
the circle $|y|=y_*$ (which encloses $0$, $y_1$ and some apparent points); $2$ and $3$, loops around $y_2$ and $y_3$
conjugated by paths that avoid all $21$ apparent points. The computation:
\begin{itemize}
\item initial data at $b=2^{-21}$ from the power series of $g$, with exact $Q_n(q)$, $P_{3,n}(q)$ and the tail bound of
Corollary~\ref{cor:size};
\item Taylor steps $c\mapsto c+h$ along the paths, $|h|\le0.2\delta$ ($\delta$ the distance to the nearest singular point
of $L_{13}$): coefficients of an approximate solution by the recurrence at a midpoint, the residual
$L_{13}\hat y-E$ computed in balls, and the error $y-\hat y$ bounded by a Cauchy majorant whose constants
$\max|A_i/A_{13}|$ on a circle are enclosed by arc subdivision using factored forms of the coefficients and certified root
balls of the leading coefficient; the step matrix is enclosed and applied as a ball product;
\item working precision $2200$ bits, truncation target $1700$ bits per step; all radii of the generator matrices stay
below $10^{-308}$ (typically $10^{-400}$).
\end{itemize}
The output is three $14\times14$ ball matrices $g_A,g_2,g_3$ whose last row is exactly $(0,\dots,0,1)$, and the
coordinate vector $v$ of $g$. As a check, $g_Av=v$ holds within the balls ($H_3$ is holomorphic inside $|y|<y_2$). The
$13\times13$ blocks are the actions on $S$.

\emph{The reflection criterion.} Let $u=(g_2-1)v\in S$.
\begin{enumerate}
\item[(a)] If the $13$ vectors $g_wu$ ($w$ words in $g_A,g_2,g_3$) are linearly independent, then $S=\operatorname{span}
\{g_wu\}$. Each $g_wu$ is a continuation of $(\gamma_2-1)H_3$, a solution of the homogeneous equation $Lf=0$, so
$S\subseteq\Sol(L)$, and loops act on $S$ through their classes in $\pi_1(\CC\setminus\{0,y_1,\dots,y_4\})$. In
particular $g_2$ acts on the affine module $\CC H_3\oplus S$ as a conjugate of the local loop at $y_2$, which is a
reflection by Proposition~\ref{prop:exponents}. Its image is $\CC u$, so $g_2|_S=1+u\otimes\phi$ for some $\phi\in S^*$.
\item[(b)] The row $\psi=$ row $i_0$ of $(g_2-1)|_S$ equals $u_{i_0}\phi$; one entry ball excludes $0$, so $\phi\ne0$
and $u_{i_0}\ne0$.
\item[(c)] If moreover the $13$ covectors $\psi g_w$ are linearly independent, $S$ is irreducible under
$\langle g_A,g_2,g_3\rangle$: for an invariant subspace $W\ne0$, either $\phi(W)\ne0$, and then
$u=(g_2x-x)/\phi(x)\in W$ for some $x\in W$, so $W\supseteq\operatorname{span}\{g_wu\}=S$; or $\phi(W)=0$, and then
$W\subseteq\bigcap_w\ker(\phi g_w)=0$.
\end{enumerate}

\begin{proposition}\label{prop:irred}
For every $q$ with $26\le q\le53$, the $13\times13$ determinants of (a) and (c) exclude $0$. Hence $S$ is irreducible
under the monodromy group of $L_{13}$.
\end{proposition}

\begin{proof}[Proof (computer-assisted, Arb)]
The words were chosen greedily on midpoints (a heuristic that does not affect rigour); at every $q$ both lists are
$\varnothing,A,A^2,A^3,A^4,3A^3,A^5,A3A^3,A^6,A^23A^3,A^7,A^33A^3,A^8$ (in the notation of the certificate files). The
determinants are evaluated in Arb at $2200$ bits from the ball entries (\file{fourterm/irred.py}); their imaginary parts
are listed in Table~\ref{tab:4T}, with relative radii at most $4\cdot10^{-30}$; the real parts are balls around $0$ of
radius at most $2\cdot10^{-219}$. Irreducibility under the subgroup $\langle g_A,g_2,g_3\rangle$ implies irreducibility under the
whole monodromy group.
\end{proof}

\begin{table}
\caption{Certificates for Theorem~\ref{thm:4T}: the least $\fm$ with $M_\fm(q)>0$, the margin $M_{14}(q)$ (midpoint of
an Arb ball of radius $\le5\cdot10^{-8}$, truncated to six decimals), $y_2$, the base point $y_*$, and the midpoints of the two $13\times13$
determinants of Proposition~\ref{prop:irred} (\file{code/thm1\_cor2/fourterm/summary\_4t.json}). The jump between
$q=37$ and $q=38$ is a change of scaled coordinates (base point and scaling), not a failure.}\label{tab:4T}
\small
\begin{tabular}{rrrrrll}
\toprule
$q$ & $\fm_{\min}$ & $M_{14}(q)$ & $y_2$ & $y_*$ & det (a) & det (c)\\
\midrule
26 & 14 & 0.009225 & 156.368 & 66 & $1.634\cdot10^{14}\,i$ & $2.307\cdot10^{89}\,i$ \\
27 & 12 & 0.082012 & 169.118 & 71 & $1.605\cdot10^{13}\,i$ & $1.657\cdot10^{91}\,i$ \\
28 & 11 & 0.152055 & 182.368 & 77 & $1.985\cdot10^{12}\,i$ & $1.121\cdot10^{93}\,i$ \\
29 & 10 & 0.219554 & 196.119 & 82 & $2.151\cdot10^{11}\,i$ & $5.668\cdot10^{94}\,i$ \\
30 & 9 & 0.284686 & 210.369 & 88 & $2.871\cdot10^{10}\,i$ & $2.739\cdot10^{96}\,i$ \\
31 & 8 & 0.347612 & 225.119 & 95 & $4.584\cdot10^{9}\,i$ & $1.253\cdot10^{98}\,i$ \\
32 & 7 & 0.408477 & 240.369 & 101 & $6.520\cdot10^{8}\,i$ & $6.072\cdot10^{99}\,i$ \\
33 & 7 & 0.467411 & 256.119 & 108 & $1.094\cdot10^{8}\,i$ & $5.079\cdot10^{101}\,i$ \\
34 & 7 & 0.524534 & 272.370 & 114 & $1.669\cdot10^{7}\,i$ & $2.784\cdot10^{103}\,i$ \\
35 & 6 & 0.579952 & 289.120 & 121 & $2.964\cdot10^{6}\,i$ & $1.637\cdot10^{105}\,i$ \\
36 & 6 & 0.633764 & 306.370 & 129 & $6.013\cdot10^{5}\,i$ & $1.007\cdot10^{107}\,i$ \\
37 & 6 & 0.686062 & 324.120 & 136 & $1.120\cdot10^{5}\,i$ & $4.311\cdot10^{108}\,i$ \\
38 & 5 & 0.736928 & 342.370 & 144 & $7.143\cdot10^{27}\,i$ & $3.168\cdot10^{92}\,i$ \\
39 & 5 & 0.786438 & 361.120 & 152 & $1.535\cdot10^{27}\,i$ & $6.101\cdot10^{93}\,i$ \\
40 & 5 & 0.834663 & 380.370 & 160 & $3.367\cdot10^{26}\,i$ & $1.076\cdot10^{95}\,i$ \\
41 & 5 & 0.881668 & 400.120 & 168 & $7.537\cdot10^{25}\,i$ & $1.747\cdot10^{96}\,i$ \\
42 & 5 & 0.927512 & 420.371 & 177 & $1.865\cdot10^{25}\,i$ & $2.758\cdot10^{97}\,i$ \\
43 & 5 & 0.972252 & 441.121 & 185 & $4.342\cdot10^{24}\,i$ & $3.843\cdot10^{98}\,i$ \\
44 & 4 & 1.015940 & 462.371 & 194 & $1.110\cdot10^{24}\,i$ & $5.232\cdot10^{99}\,i$ \\
45 & 4 & 1.058624 & 484.121 & 203 & $2.884\cdot10^{23}\,i$ & $9.447\cdot10^{100}\,i$ \\
46 & 4 & 1.100350 & 506.371 & 213 & $8.133\cdot10^{22}\,i$ & $2.221\cdot10^{102}\,i$ \\
47 & 4 & 1.141158 & 529.121 & 222 & $2.178\cdot10^{22}\,i$ & $4.203\cdot10^{103}\,i$ \\
48 & 4 & 1.181090 & 552.371 & 232 & $6.302\cdot10^{21}\,i$ & $8.202\cdot10^{104}\,i$ \\
49 & 4 & 1.220180 & 576.121 & 242 & $1.847\cdot10^{21}\,i$ & $1.468\cdot10^{106}\,i$ \\
50 & 4 & 1.258466 & 600.371 & 252 & $5.482\cdot10^{20}\,i$ & $2.422\cdot10^{107}\,i$ \\
51 & 4 & 1.295978 & 625.121 & 263 & $1.740\cdot10^{20}\,i$ & $4.090\cdot10^{108}\,i$ \\
52 & 4 & 1.332747 & 650.371 & 273 & $5.295\cdot10^{19}\,i$ & $5.802\cdot10^{109}\,i$ \\
53 & 4 & 1.368803 & 676.122 & 284 & $1.717\cdot10^{19}\,i$ & $8.436\cdot10^{110}\,i$ \\
\bottomrule
\end{tabular}
\end{table}

\begin{proof}[Proof of Proposition~\ref{prop:indep}, Theorem~\ref{thm:4T} and Corollary~\ref{cor:F1}]
By Lemma~\ref{lem:nonzero}, $U_a$ is a nonzero monodromy-invariant subspace of $S$, so $U_a=S$ by
Proposition~\ref{prop:irred}; then Lemma~\ref{lem:indep} with $u=13$ gives Proposition~\ref{prop:indep}, and the
reduction above gives Theorem~\ref{thm:4T}. Corollary~\ref{cor:F1} is the case $26\le q\le32$ of Theorem~\ref{thm:4T}:
if $\theta_q=u/v$ with $u,v\in\ZZ$, $v\ne0$, then $(r_0,a)=(-u,0,0,v)$ is a nonzero relation.
\end{proof}

\begin{remark}[Independent certificate for Corollary~\ref{cor:F1}]\label{rem:F1cert}
Corollary~\ref{cor:F1} was first proved by a different final step: for $26\le q\le32$ the $14$ continuations of $H_3$
along the words $\varnothing,2,3,A2,A3,A^22,A^23,A^32,A^33,A^42,A^43,A^52,A^62,A^72$ have a nonzero $14\times14$ Arb
determinant (\file{cert\_det.py}), for example
\[
\begin{gathered}
\relax[257906186217.156323475542228360\pm3.7\cdot10^{-19}]\,i\quad(q=26),\\
[992625.090850805422458198252898\pm3.9\cdot10^{-25}]\,i\quad(q=32).
\end{gathered}
\]
So
$\dim\operatorname{span}\{\gamma H_3\}=14$, and the elementary analogue of Lemma~\ref{lem:indep} gives the independence
of $1,H_3,\dots,\vartheta^{12}H_3$. (The passage from coordinate vectors to germs uses $E\ne0$.) Independent reruns at
$q=26$ and $q=32$ with initial data that do not use Lemma~\ref{lem:contact} (base $b=2^{-22}$, the crude bound
$(2n+1)2^{20n}|\mu_3|/(q-1)$) and a different discretisation reproduce these determinants to all $30$ printed digits,
and give the same irreducibility certificates.
\end{remark}

\begin{remark}
Smaller holonomy budgets suffice for larger $q$ ($\fm_{\min}$ in Table~\ref{tab:4T}), but the proof uses $\fm=14$ at
every $q$. No case analysis on the relation vector and no input from linear forms in logarithms is needed: only $a\ne0$
is used.
\end{remark}

% ===========================================================================
\section{Widening CDT's range: proof of Theorem~\ref{thm:Cprime}}\label{sec:Cprime}
% ===========================================================================

\subsection{CDT's construction}\label{sec:CDTconstruction}
We recall \cite[\S14]{CDT}. For an odd integer $a$ with $|a|\ge3$ let $A(a,x)=(1-2ax+x^2)^{-1/2}=\sum_ku_k(a)x^k$ and
$H(a,x)=A(a,x)\int_0^xA(a,t)\,dt=\sum_kv_k(a)x^k$, so that $v_k(a)/u_k(a)\to\frac12\log\frac{a+1}{a-1}$. For distinct
odd $a=2m+1$, $b=2n+1$ put $\lambda_a=\frac12\ell_m$, $\lambda_b=\frac12\ell_n$, and suppose
$r_0+r_a\lambda_a+r_b\lambda_b+r_{ab}\lambda_a\lambda_b=0$ with $r\in\ZZ^4$, $(r_a,r_b,r_{ab})\ne0$. CDT form
\[
P(x)=\sum_k\bigl(r_0u_k(a)u_k(b)+r_av_k(a)u_k(b)+r_bu_k(a)v_k(b)+r_{ab}v_k(a)v_k(b)\bigr)x^k
\]
(a combination of Hadamard products) and symmetrise:
$G(y)=P(x)+P(x/(x-1))$ with $y=x^2/(x-1)$ \cite[Def.~14.1.9]{CDT}. By \cite[Prop.~14.1.11]{CDT}, $G$ has denominator
type $[1,\dots,2k]^2$, satisfies a differential equation $\mathcal L_{a,b}$ whose singular points are $0,4,\infty$ and
the four points $Z_4(a,b)$, the $y$-images of the four products $x=(a\pm\sqrt{a^2-1})(b\pm\sqrt{b^2-1})$ (plus apparent
singularities at the roots of a polynomial $R_{10}$), and $G$ is overconvergent at the image $\alpha$ of the smallest
product. With the seven ``pure'' functions $B_1,\dots,B_7$ of \cite[\S10]{CDT} and ten functions built from $G$
($G,G',G'',G'''$ and six primitives), \cite[Lemma~14.3.1]{CDT} gives $\fm=17$ functions, linearly independent over
$\QQ(y)$, with
\[
\begin{gathered}
\sigma=(0,2,2,4,\dots,4),\qquad e=(0,0,1;0,0,0,0,0,0;1,1,1,1,1,1,1,1),\\
\tau=\frac{1136}{289}+\frac{78419}{242760}=\frac{1032659}{242760}=4.2538\ldots
\end{gathered}
\]
\cite[(14.5.2)--(14.5.4)]{CDT}. The order of the functions is CDT's:
$B_1,B_2,B_3$; $B_4,B_5,G,G',G'',G'''$; $B_6,B_7$ and the six primitives.
The map is $\varphi=h\circ\psi$, with
\[
h(\mathsf q)=-256\,\mathsf q\prod_{k\ge1}(1+\mathsf q^k)^{24}=-256\Bigl(\frac{\eta(2\tau)}{\eta(\tau)}\Bigr)^{24},
\]
the hauptmodul of $\Gamma_0(2)$ in the cusp coordinate $\mathsf q=e^{2\pi i\tau}$, and $\psi:\DD\to\DD$ holomorphic
with $\psi^{-1}(0)=\{0\}$; admissibility of $\varphi$ is checked through \cite[Cor.~9.0.19]{CDT}. The constant
$10^{-6}$ of \cite[Thm.~C]{CDT} enters only through \cite[Lemma~14.2.4]{CDT} (which needs $|a/b\mp1|<\frac12$),
\cite[Lemma~14.2.6]{CDT} (which needs $0<|a/b-1|<10^{-3}$), and the location of the singularities, \cite[\S14.4]{CDT},
which determines the admissible maps. We change only the map.

\subsection{Location of the singular values}

\begin{lemma}\label{lem:S}
Let $0<\varepsilon\le10^{-4}$ and let $m\ne n$ be integers in $\ZZ\setminus\{-1,0\}$ with $|m/n-1|<\varepsilon$; put
$a=2m+1$, $b=2n+1$. Then the three genuine singular values (the points of $Z_4(a,b)$ other than $\alpha$) are real with
$|y|\ge Y_1:=1/\varepsilon-6$, and $\alpha$ is real with $|\alpha|\le Y_0:=\varepsilon^4/16$.
\end{lemma}

\begin{proof}
Let $A=|a|$, $B=|b|$, $u=a+\operatorname{sgn}(a)\sqrt{a^2-1}$, $v$ likewise; the products are $x\in\{uv,u/v,v/u,1/(uv)\}$,
all real since $A,B\ge3$, and $y=x^2/(x-1)=(x-1)+2+1/(x-1)$. By the mean value theorem for $\operatorname{arccosh}$,
$\delta:=|\log|u|-\log|v||\le|A-B|/\sqrt{\min(A,B)^2-1}$. For positive pairs $|A-B|/\min(A,B)<\varepsilon/(1-\varepsilon)$,
for negative pairs $\le2\varepsilon|n|/(2|n|(1-\varepsilon)-1)\le\varepsilon/(1-1.5\varepsilon)$, using $|n|\ge1/\varepsilon$;
and $\min(A,B)\ge2/\varepsilon-3$, so the square-root factor is at most $1+\varepsilon^2/7$. Hence
$\delta\le\varepsilon(1+1.6\varepsilon)$, and $x\in\{u/v,v/u\}$ has $|x-1|\le e^\delta-1\le\varepsilon(1+3\varepsilon)$, so
$|y|\ge(1-3\varepsilon)/\varepsilon-2-2\varepsilon>1/\varepsilon-6$. For $x=uv$, $|x|\ge(2\min(A,B)-1)^2$ gives $y>1/\varepsilon$;
for $x=1/(uv)$, $|\alpha|\le2/|uv|^2<\varepsilon^4/16$.
\end{proof}

The constant $6$ cannot be dropped entirely: CDT's ``$|y|>1/\varepsilon$'' fails by $O(1)$. At
$(m,n)=(-10^6,-10^6-1)$ the smallest genuine singular value has $|y|=10^6-1.50$; at $\varepsilon=3\cdot10^{-6}$ the worst
case is $1/\varepsilon-1.83$, at $(m,n)=(-333333,-333334)$. This does not affect \cite[Thm.~C]{CDT}, whose map has room.

\subsection{Admissibility of one map for the whole range}
Let $\varphi=h\circ\psi$ with $\psi$ entire, $\psi^{-1}(0)=\{0\}$, and suppose
\begin{enumerate}
\item[(a)] $|\psi|<1$ on $\TT$;
\item[(b)] $|\varphi|>2Y_0$ on $\TT$;
\item[(c$'$)] $|\varphi|<Y_1$ on $\TT$.
\end{enumerate}
By (a) and the maximum principle $\psi(\overline\DD)\subset\DD$, so $\varphi$ is holomorphic on a neighbourhood of
$\overline\DD$. By (c$'$), $\varphi(\DD)\subset D(0,Y_1)$, which by Lemma~\ref{lem:S} misses every genuine singular value
of every pair in the range. Since $\varphi$ vanishes only at $0$, simply, and $|\varphi|>2Y_0$ on $\TT$, the argument
principle shows that every point of $U=D(0,2Y_0)$ has exactly one preimage in $\DD$; this is the univalent leaf over
$U\ni\alpha$ required by \cite[Cor.~9.0.19]{CDT}, with $Y^0=\{\alpha\}$ and $Y^1$ the three genuine singular values.
(Cor.~9.0.19 leaves $U$ free, subject to being a contractible neighbourhood of $0$ in $\CC\setminus\{4\}$ containing
$\alpha$ on which $G$ is holomorphic; $G$ is holomorphic on its principal disc, which contains $D(0,2Y_0)$.) So the
$17$ pullbacks $f_i\circ\varphi$ are meromorphic on $\DD$, for every pair with $|m/n-1|<\varepsilon$ at once.

\subsection{The certificate}\label{sec:certC}
We take $\psi(z)=z\exp(\sum_{k=0}^{96}c_kz^k)$ with dyadic coefficients $c_k\in2^{-40}\ZZ$ (the certified object; stored
in \file{code/thm3/runs/cert3\_p4\_C\_eps3.0e-06\_eps3.0e-06\_l3.json}), found by an admissibility-aware L-BFGS
optimisation of $L-T/17$ (started from lune-shaped maps and from CDT's map). Then $L=\log|\varphi'(0)|=\log256+c_0$ exactly. We apply
Theorem~\ref{thm:CDT602} with the four maps $\varphi_k(z)=\varphi(r_kz)$, where
\[
(r_0,r_1,r_2,r_3)=\Bigl(\frac1{2048},\frac{99}{256},\frac{215}{256},1\Bigr),\quad
(\vartheta_1,\vartheta_2,\vartheta_3)=\Bigl(\frac{61}{200},\frac{231}{200},\frac{1109}{500}\Bigr);
\]
the $\varphi_k$ are restrictions of $\varphi$, so they are admissible.

\emph{Certification} (\file{cert3.py}, \file{cert3d.py}). The circle is cut into dyadic intervals, adaptively. On each
interval, an Arb value of $\log\varphi$ at the midpoint (via $\eta$-quotients) and a rigorous bound for
$|d\log\varphi/d\vartheta|$, namely $(1+24\sum_kk\rho^k/(1-\rho^k))\cdot\sup|1+zg'(z)|$ with $\rho\ge\sup|\psi|$ on the
interval and $g=\sum c_kz^k$, give an upper step function for $\log|\varphi|$, hence an exact rational upper bound for
the rearrangement integral of the mixture, and decide (a), (b), (c$'$). The results at $\varepsilon=3\cdot10^{-6}$:
\begin{itemize}
\item (a), (b), (c$'$) hold; $\max_\TT\log|\varphi|\le13801311154345/2^{40}=12.55222<\log Y_1=12.71688$;
\item $L=4.7983738127618231\ldots$, mixture integral $\le9.08327$, penalty
$\frac1{17}\sum\vartheta_k^2\log(r_k/r_{k-1})=2.51411\ldots/17$, so that the numerator is
$[9.23115527305872\pm1.6\cdot10^{-15}]$;
\item hence $17\le\text{numerator}/(L-\tau)\le16.95200$ would be needed: contradiction, with margin
$L-\text{numerator}/17-\tau\ge0.0015378$.
\end{itemize}
The margin is thin ($0.048$ in the function count) but certified, and all slack of the step functions is in the safe
direction. Certificates for $\varepsilon=2\cdot10^{-6}$ and $2.5\cdot10^{-6}$ (trivial partition; margins $\ge0.01138$
and $\ge0.00131$) and a certificate for CDT's own map at $10^{-6}$ (margin $\ge0.0308$) serve as cross-checks; only the
$3\cdot10^{-6}$ certificate is needed, since Lemma~\ref{lem:S} at $\varepsilon=3\cdot10^{-6}$ covers every pair below it.
An independent non-rigorous sampling with a $\mathsf q$-product instead of $\eta$ gives numerator $\approx9.2297$.

\subsection{The \texorpdfstring{$\varepsilon$}{epsilon}-dependent steps of CDT's chain}\label{sec:phi6}
Lemma~14.2.4 of \cite{CDT} needs $|a/b\mp1|<\frac12$, and Lemma~14.2.6 needs $0<|a/b-1|<10^{-3}$; here
$|a/b-1|<\varepsilon(1+2\varepsilon)\le10^{-5}$. The proof of Lemma~14.2.6 uses only $b-a=2(n-m)\in\ZZ\setminus\{0\}$, so
$|m-n|\ge2$ changes nothing. \cite[Lemma~14.3.1]{CDT} needs a $10\times10$ determinant $\det(a,b)$ to be nonzero. We
recomputed it exactly from the operator of \cite[Prop.~14.1.11]{CDT} and the reduction \cite[(14.3.2)--(14.3.5)]{CDT},
with numerators $S_k$ of the shape \cite[(14.3.6)]{CDT} (one for each basis vector of the unknown $Q$), and proved (Proposition~\ref{prop:detid} below)
\[
\det(a,b)(a-b)^{86}=\kappa\,a^6b^6(a+b)^{24}Q_1^4Q_2^4Q_3^4Q_4^2Q_5^2\,\Phi_4\Phi_6\Phi_{14}^3\Phi_{26}\Phi_{79}^3,
\]
with $\kappa=645700815/37778931862957161709568$, $Q_1=4a^2-4ab+1$, $Q_2=4b^2-4ab+1$, $Q_3=4a^2-4ab+4b^2-3$,
$Q_4=16a^2-40ab+16b^2+9$, $Q_5=80a^2-128ab+80b^2-45$, and $\Phi_d$ irreducible of $a$-degree $d$. The factors are handled as follows.
\begin{itemize}
\item $a$, $b$, $a\pm b$ are nonzero, and the $Q_i$ are odd for odd $a,b$.
\item $\Phi_4=80a^4-88a^3b+24a^2b^2-399a^2-88ab^3+786ab+80b^4-399b^2+9$ is odd for odd $a,b$, hence nonzero. (CDT's
Lemma~14.2.6 asserts $\Res(R_{10},y(4-y))\ne0$, but its proof treats only $\Phi_{14}$ and $\Phi_{79}$.)
\item For $\Phi_{14}$ and $\Phi_{79}$ we use the argument in the proof of CDT's Lemma~14.2.6, which uses only
$|b-a|\ge1$ and $|a/b-1|\le10^{-3}$; a larger $|b-a|$ only shrinks the error terms.
\item $\Phi_{26}$: Lemma~14.2.4 gives $\Res(R_4,R_{10})\ne0$ by irreducibility and division, for $|a/b-1|<\frac12$.
\item $\Phi_6$ divides none of $\Delta_y(R_{10})$, $\Res_y(R_{10},R_4)$, $\Res_y(R_{10},y(y-4))$, contrary to the text of
\cite{CDT}, so neither lemma covers it. Explicitly,
\[
\begin{aligned}
\Phi_6(a,b)={}&59290a^6b^2-185493a^5b^3+249865a^4b^4-185493a^3b^5+59290a^2b^6\\
&+8491a^5b-65541a^4b^2+119686a^3b^3-65541a^2b^4+8491ab^5\\
&-6090a^4+15225a^3b-21000a^2b^2+15225ab^3-6090b^4\\
&+975a^2-4260ab+975b^2+1755,
\end{aligned}
\]
and $\Phi_6(a,b)\equiv16\pmod{32}$ for all odd $a,b$ (check over the $16\times16$ odd residue classes mod $32$). So
$\Phi_6(a,b)\ne0$ at every odd pair. (Also $|\Phi_6(a,b)|/|a|^8\ge1864$ whenever $a\ne b$ are integers with
$|a/b-1|<1/20$, so $\Phi_6$ has no integer zeros near the diagonal.)
\end{itemize}
Hence, given CDT's set-up, \cite[Lemma~14.3.1]{CDT} holds for all odd $a\ne b$ with $|a/b-1|<10^{-3}$, with no extra
hypothesis. The set of irreducible factors agrees with CDT's printed list; the multiplicities differ, which does not
matter. A misprint in \cite{CDT}: the factor ``$1+4b^2-2ab$'', printed in both (14.2.2) and (14.2.3), should read
$1+4b^2-4ab$.

\begin{proposition}\label{prop:detid}
Given the shape \cite[(14.2.1)]{CDT} of the operator of \cite[Prop.~14.1.11]{CDT} and the reduction
\cite[(14.3.2)--(14.3.5)]{CDT}, the displayed identity for $\det(a,b)(a-b)^{86}$ holds in $\QQ(a,b)$, hence at every
pair of integers $a\ne b$.
\end{proposition}

\begin{proof}[Proof (computer-assisted, exact)]
\emph{The operator.} CDT print only $R_4$, $R_{10}$ and $c_3/c_4$ of their operator $\mathcal L=\sum_{i\le4}c_iD^i$. We
interpolate a candidate $\tilde{\mathcal L}\in\QQ[a,b,y][D]$ with $c_4=(y-4)^2y^3R_4R_{10}$ and verify it: the
coefficients of $y^k$, $k<63$, of $\tilde{\mathcal L}(G_A)$ vanish identically in $\QQ[a,b]$, where
$G_A=\sum_nP_n(a)P_n(b)p_n(y)$ ($P_n$ Legendre polynomials, $p_n(x+w)=x^n+w^n$ for $w=x/(x-1)$), and these $63$
conditions on the $63$ unknowns of CDT's ansatz have rank $63$ (already at $(a,b)=(7,5)$). So the ansatz has a unique
solution over $\QQ(a,b)$, and CDT's operator is $\tilde{\mathcal L}$; also $c_3/c_4$ agrees symbolically with CDT's
printed formula (\file{f1\_Lsym.py}).

\emph{An a priori degree bound.} With $\tilde{\mathcal L}$, the $10$ numerators $S_k=y^4R_{10}^4\sum_ia_iy^i$ obtained
from $b_3=(y-4)^2R_4y^k/R_{10}$, $k=0,\dots,9$ (CDT's (14.3.6) is the case of a general $Q=\sum_kq_ky^k$, with one
numerator $S_{42}$), are explicit in $\QQ[a,b,y]$. In the variables $t=a-b$ and $b$, the leading coefficient of $R_{10}$ is $-32t^4$, and the
remainders of $S_k$ modulo $R_{10}$ give $\det(a,b)=\det M'(t,b)/t^{200}$ with $M'$ a polynomial matrix. Summing the
largest entry degrees column by column gives $\deg_b\det M'\le690$ (and $\deg_t\det M'\le685$), using the entries only
(\file{f1\_detsym.py}).

\emph{The identity.} It is equivalent to $\det M'(t,b)-\kappa\,t^{114}\Pi(t+b,b)=0$, where $\kappa\Pi$ is the right
side; this polynomial has $b$-degree at most $\max(690,\deg\Pi)=690$. For each of the $692$ integers
$b_0=-346,\dots,345$ we computed $\det M'(t,b_0)\in\QQ[t]$ exactly (fraction-free elimination, every division checked
exact) and found it equal to $\kappa t^{114}\Pi(t+b_0,b_0)$ (\file{f1\_identity.py}, about $35$ minutes). A polynomial of
$b$-degree $\le690$ that vanishes on $692$ lines $b=b_0$ is zero.
\end{proof}

An earlier version of this computation interpolated $\det(a,b_0)(a-b_0)^{86}$ on ten lines from $420$ values each, which
gives the right answer (it is a consequence of Proposition~\ref{prop:detid}) but, by the bound above, did not by itself
prove it (\file{code/thm3/f1/}).

\begin{proof}[Proof of Theorem~\ref{thm:Cprime}]
A relation for a pair with $|m/n-1|<3\cdot10^{-6}$ gives, by \cite[Prop.~14.1.11, Lemma~14.3.1]{CDT} and
Section~\ref{sec:phi6}, $17$ functions as in Section~\ref{sec:CDTconstruction}, linearly independent over $\QQ(y)$ and
holonomic. The map of Section~\ref{sec:certC} is admissible for them, and Theorem~\ref{thm:CDT602} gives
$17\le16.952$, a contradiction.
\end{proof}

% ===========================================================================
\section{Pooling pairs: proof of Theorem~\ref{thm:F2}}\label{sec:F2}
% ===========================================================================

Suppose there are $K$ pairwise inequivalent exceptional classes with representatives $(m_k,n_k)$ in the range. Each gives
a $G$-function $G_k$ and its ten associated functions (the \emph{$G$-family} of the pair); the pure functions
$B_1,\dots,B_7$ do not depend on the pair. The idea is to put all $7+10K$ functions into one holonomy bound.

\subsection{Joint independence}

\begin{lemma}\label{lem:D}
Assume \cite[Lemma~14.2.4]{CDT}. Let $(a,b)$, $(c,d)$ be integer pairs with $|a|,|b|,|c|,|d|\ge2$, $a\ne b$, $c\ne d$,
$|a/b-1|<\frac12$, $|c/d-1|<\frac12$. If $Z_4(a,b)\cap Z_4(c,d)\ne\varnothing$, then $\{c,d\}=\{a,b\}$ or
$\{c,d\}=\{-a,-b\}$. Moreover $0,4\notin Z_4(a,b)$.
\end{lemma}

\begin{proof}
$Z_4(a,b)$ is the root set of the quartic $R_4(a,b;y)$ of \cite[(14.2.1)]{CDT}, which is irreducible over $\QQ$ by
\cite[Lemma~14.2.4]{CDT}, with constant term $1$ and leading coefficient $4(a-b)^2$. Two irreducible quartics with a
common root and constant term $1$ are equal. With $\sigma=a^2+b^2$ and $\pi=ab$ the coefficients of $y^4$ and $y^3$ are
$c_4=4\sigma-8\pi$ and $c_3=32\pi^2-16\pi\sigma+16\pi-8\sigma$, and $c_3=-2c_4-4\pi c_4$ identically. So $R_4$ determines
$\pi$ and $\sigma$, hence $(a+b)^2$ and $(a-b)^2$, hence $\{a,b\}$ up to a common sign. Finally $R_4(0)=1$ and
$R_4(4)=(16a^2-40ab+16b^2+9)^2$, which is odd (\file{code/thm4/lemmaD.py}).
\end{proof}

\begin{lemma}\label{lem:H}
Every branch of every function in the $G$-family of $(a,b)$ is holomorphic on $\CC\setminus(\{0,4\}\cup Z_4(a,b))$, and
every branch of $B_1,\dots,B_7$ is holomorphic on $\CC\setminus\{0,4\}$.
\end{lemma}

\begin{proof}[Proof sketch]
$P(x)$ is a combination of Hadamard products $f\star g$ with $f\in\{A(a,\cdot),H(a,\cdot)\}$ and
$g\in\{A(b,\cdot),H(b,\cdot)\}$, each factor continuable along every path avoiding its two singular points. By
Hadamard's multiplication theorem in its global form (deform the contour in $f\star g(x)=\frac1{2\pi i}\oint
f(t)g(x/t)\,dt/t$; \cite{Had99} for the star-domain form, and \cite[Thm.~A]{LSS20} (preprint) for the path-wise form:
if $f$ continues along every path avoiding a closed discrete set $S$ and $g$ along every path avoiding $T$, then
$f\star g$ continues along every path avoiding $\{0\}\cup S\cdot T$), $f\star g$ continues along every path avoiding $0$ and the
four products. In $y$,
the two roots of $x^2-yx+y=0$ are distinct off $y\in\{0,4\}$ and avoid $\{0\}\cup\{\text{products}\}\cup\{p/(p-1)\}$
exactly off $\{0\}\cup Z_4$, so $G=P(x)+P(x/(x-1))$ continues along every path in $\CC\setminus(\{0,4\}\cup Z_4)$;
derivatives and primitives inherit this. The $B_i$ are functions on $\mathbb P^1\setminus\{0,4,\infty\}$
\cite[\S10]{CDT}.
\end{proof}

\begin{lemma}[cross-pair independence]\label{lem:X}
Make the assumptions of Theorem~\ref{thm:Cprime}. Let $(a_k,b_k)$, $k=1,\dots,K$, be odd pairs with
$0<|a_k/b_k-1|<10^{-3}$, pairwise inequivalent ($\{a_j,b_j\}\ne\pm\{a_k,b_k\}$ for $j\ne k$), each with a relation
vector as above. Then $B_1,\dots,B_7$ and the $K$ $G$-families are linearly independent over $\CC(y)$.
\end{lemma}

\begin{proof}[Proof sketch]
Suppose $\sum_iS_iB_i+\sum_k\sum_lR_{kl}F_{kl}=0$ with $S_i,R_{kl}\in\CC(y)$, $F_{kl}$ the $G$-family of pair $k$. Fix $k$.
CDT's proof of Lemma~14.3.1 (following Lemma~12.1.1) uses a loop $\gamma_k$, a conjugate of a small loop around a point
$s\in Z_4(a_k,b_k)$, with $G_k\ne\gamma_kG_k$. By Lemma~\ref{lem:D}, $s\notin\{0,4\}\cup Z_4(a_j,b_j)$ for $j\ne k$, so by
Lemma~\ref{lem:H} every branch of every $F_{jl}$ ($j\ne k$) and of every $B_i$ is holomorphic at $s$ and fixed by
$\gamma_k$. Applying $1-\gamma_k$ leaves $\sum_lR_{kl}(1-\gamma_k)F_{kl}=0$, which is exactly the relation reached in
CDT's single-pair proof (the varied primitives $(1-\gamma_k)\int y^e(G_k-T)$ are primitives of $y^e(1-\gamma_k)G_k$,
and CDT's argument does not depend on the choice of primitive). CDT's argument, with the determinant of
Section~\ref{sec:phi6}, gives $R_{kl}=0$. Doing this for every $k$ leaves $\sum S_iB_i=0$, so $S_i=0$
\cite[Lemma~12.0.1]{CDT}.
\end{proof}

Pairs sharing an entry, $(a,b)$ and $(a,c)$, are inequivalent, so Lemma~\ref{lem:X} covers them. As a falsification
test (not used in the proof), the exact rank modulo two primes $p>2008$ of the union of the series, truncated at
$y^{1000}$, with polynomial coefficients of degree $\le D$, is full for $30$ instances ($K=1,2,3$, including shared
entries, pairs in the range and degenerate relation vectors), while duplicated pairs give exactly the expected nullity
$10(D+1)$ (\file{code/thm4/xrank.py}, \file{runs/xrank.out}).

\subsection{The pooled holonomy inequality}
By Lemma~\ref{lem:X} we have $\fm=7+10K$ functions, holonomic and linearly independent over $\QQ(y)$, with
$\sigma=(0,2,2,4,4,4,4\mid4^{10K})$ and $e=(0,0,1,0,0,1,1\mid(0^4,1^6)^K)$ (now in the order $B_1,\dots,B_7$, then for each
pair $G,G',G'',G'''$ and the six primitives; for $K=1$ this is CDT's list reordered within the block $\sigma_i=4$, which
changes neither $\tau(b)$ nor $\tau^\sharp(e)$). Evaluating $\tau^\sharp$ of
\cite[(6.0.6)]{CDT} at the explicit point $\xi=169/100$ (an upper bound for the minimum, hence a valid $\tau$) gives
\[
\tau_2=\frac{1106977681}{262702440}=4.2138081\ldots,\qquad\tau_3=\frac{29601442477}{7082066992}=4.1797744\ldots
\]
for $K=2$ and $K=3$
(\file{tau\_pool.py}; CDT's $\tau$ for $K=1$ is reproduced by the same code).

For the maps of this section the disc condition (c$'$) fails, and we use instead:

\begin{lemma}[ray lemma]\label{lem:ray}
Let $\varphi$ be holomorphic on a neighbourhood of $\overline\DD$ and $\Gamma=(-\infty,-Y_1]\cup[Y_1,\infty)$. If
$\varphi(\TT)\cap\Gamma=\varnothing$, then $\varphi(\overline\DD)\cap\Gamma=\varnothing$.
\end{lemma}

\begin{proof}
For $y_0\in\Gamma$ the number of zeros of $\varphi-y_0$ in $\DD$ is the winding number of $\varphi(\TT)$ about $y_0$. It
is constant on each half-line of $\Gamma$ (connected subsets of $\CC\setminus\varphi(\TT)$), and $0$ for
$|y_0|>\max_\TT|\varphi|$.
\end{proof}

With Lemma~\ref{lem:S} (real singular values), (a), (b) and ``$\varphi(\TT)\cap\Gamma=\varnothing$'' make $\varphi$
admissible for all pairs in the range simultaneously. The certificate decides $\varphi(s)\notin\Gamma$ on each interval
either by $|\varphi|<Y_1$ or by a lower bound for $|\sin\arg\varphi|$. The results (\file{cert\_pool.py}, trivial
partition, $\delta_{\max}=0.002$):
\begin{center}
\small
\begin{tabular}{llllll}
\toprule
range $\varepsilon$ & $K$ & $\fm$ & $L$ & numerator $\le$ & certified margin $L-R^+/\fm-\tau$\\
\midrule
$10^{-5}$ & 2 & 27 & $4.7056374$ & $8.795131$ & $\ge0.16608$\\
$3\cdot10^{-5}$ & 2 & 27 & $4.5683356$ & $8.064226$ & $\ge0.05585$\\
$5\cdot10^{-5}$ & 2 & 27 & $4.5159123$ & $7.819914$ & $\ge0.01247$\\
$10^{-4}$ & 3 & 37 & $4.4006730$ & $7.326966$ & $\ge0.02287$\\
\bottomrule
\end{tabular}
\end{center}
The $5\cdot10^{-5}$ map was re-optimised for $L-T/27$ (\file{reopt.py}); the $10^{-5}$ map is an earlier optimised map
dilated by $r=1-2^{-14}$ (the undilated map crosses the ray $[Y_1,\infty)$ near $1/\varepsilon$, which the certificate
detects); the other two are stored optimised maps. Only the rows $5\cdot10^{-5}$ and $10^{-4}$ are needed. Lemma~\ref{lem:S}
is proved for $\varepsilon\le10^{-4}$; at $\varepsilon=5\cdot10^{-5}$ the map has little room ($\varphi(\TT)$ crosses the
negative axis at $\approx-19988$, inside $Y_1=19994$), and an exhaustive high-precision check of Lemma~\ref{lem:S} on a
boundary window found the closest genuine value at $Y_1+4.25$.

\begin{proof}[Proof of Theorem~\ref{thm:F2}]
Suppose there were $K=2$ (resp.\ $3$) inequivalent exceptional classes with representatives in $|m/n-1|<5\cdot10^{-5}$
(resp.\ $10^{-4}$). CDT's construction gives a $G$-function for each, and Lemma~\ref{lem:X} (applicable since
$|a/b-1|<10^{-3}$) gives $7+10K$ independent functions. The certified map is admissible for all of them, and
Theorem~\ref{thm:CDT602} gives $27\le25.885$ (resp.\ $37\le33.169$), a contradiction. This proves (i) and (ii); (iii) is
Theorem~\ref{thm:Cprime}. For (iv): a consecutive pair $(N,N+1)$ in these ranges has $1/(N+2)<10^{-4}$, so
$q=N+1\ge10^4>54$, and Hata's theorem gives the four-term independence.
\end{proof}

% ===========================================================================
\section{What is computer-assisted, and how to reproduce it}\label{sec:computer}
% ===========================================================================

The folder \file{code/} of the accompanying repository \url{https://github.com/telaviv/log-products-irrationality}
contains every script and certificate output used above, in one directory per
result; \file{code/README.md} maps each number of this paper to a script and an output file. Software: Python 3.12,
python-flint 0.9.0 (FLINT/Arb ball arithmetic), SymPy 1.14, mpmath 1.3, NumPy and SciPy (optimisation only). The
computer-assisted steps are:
\begin{center}
\small
\begin{tabular}{p{4.3cm}p{1.4cm}p{3.6cm}p{1.6cm}}
\toprule
step & kind & script (folder) & time\\
\midrule
Relations for $c_{n,l}$; order-4 recurrence; char.\ polynomial (Prop.~\ref{prop:I1}) & exact & \file{i1\_*.py}
(\file{thm1\_cor2}) & minutes\\
Contact orders, sanity checks (Lemma~\ref{lem:contact}) & exact & \file{i1\_contact.py} & 10 s\\
$L_{13}$ and I1$'$ per $q$ (Prop.~\ref{prop:I1prime}) & exact & \file{findop.py}, \file{i1\_rec25.py},
\file{fourterm/mkrec.py} & seconds per $q$\\
Budget $M_{14}(q)$ & Arb & \file{fourterm/budget\_4t.py} & seconds\\
Local exponents, Casoratian, $p_0$ roots (Prop.~\ref{prop:exponents}, Lemma~\ref{lem:nonzero}) & exact &
\file{exponents.py} (\file{fourterm}) & 15 s\\
Certified continuation, $g_A,g_2,g_3$ & Arb & \file{mono13.py} & $\approx40$ min per $q$\\
Irreducibility (Prop.~\ref{prop:irred}) & Arb & \file{fourterm/irred.py} & 2 s per $q$\\
Rank-14 cross-check (Remark~\ref{rem:F1cert}) & Arb & \file{cert\_det.py} & seconds\\
Map certificate, Theorem~\ref{thm:Cprime} & Arb + exact & \file{cert3.py}, \file{cert3d.py} (\file{thm3}) &
$\approx8$ min\\
$\det(a,b)$ identity and $\Phi_6$ (Prop.~\ref{prop:detid}) & exact & \file{f1\_Lsym.py}, \file{f1\_detsym.py},
\file{f1\_identity.py}, \file{f1\_phi6\_nonvanish.py} (\file{thm3/f1}) & 35 min\\
New pairs \eqref{eq:newpairs} & exact & \file{f7\_newpairs.py} & seconds\\
Lemma~\ref{lem:D}; pooled $\tau$; falsifier & exact & \file{lemmaD.py}, \file{tau\_pool.py}, \file{xrank.py}
(\file{thm4}) & minutes\\
Pooled map certificates (Theorem~\ref{thm:F2}) & Arb + exact & \file{cert\_pool.py} & $\approx4$ min each\\
\bottomrule
\end{tabular}
\end{center}
The map optimisations (\file{lib\_opt.py}, \file{lib\_range.py}, \file{reopt.py}) and the falsifier are not part of any
proof: they produce candidates that the certificates then check. The referees (Section~\ref{sec:provenance}) re-ran the
main certificates bit for bit on separate copies (among them the budget and irreducibility certificates of
Theorem~\ref{thm:4T}, the rank certificate at $q=29$, the $3\cdot10^{-6}$ map certificate and the $5\cdot10^{-5}$ pooled
certificate), and re-ran the continuation step independently, with different initial data and discretisation, at
$q=26,32$ (Remark~\ref{rem:F1cert}; outputs in \file{code/thm1\_cor2/runs/referee\_rerun/}) and at $q=40,53$ (the gate
referee's reruns of Theorem~\ref{thm:4T}, which reproduce both determinants of Proposition~\ref{prop:irred} to all $30$
printed digits; their outputs are not included in \file{code/}).

What the computer establishes is, in each case, a finite statement: an identity of explicit polynomials, a determinant
ball that excludes $0$, or an exact rational upper bound for an integral. What it does not establish: the hand arguments
(Lemmas~\ref{lem:den}, \ref{lem:indep}, \ref{lem:nonzero}, the reflection criterion, Lemmas~\ref{lem:S}, \ref{lem:D},
\ref{lem:ray}, the reductions), the theorems of \cite{CDT} we use, and the correctness of python-flint's Arb routines
(in particular \texttt{acb.modular\_eta} and matrix determinants), which we take on trust. Arb \cite{Arb} is now part of FLINT \cite{FLINT}; python-flint
\cite{pyflint} 0.9.0, as used, is linked against FLINT 3.6.0.

% ===========================================================================
\section{Limits}\label{sec:limits}
% ===========================================================================

\begin{itemize}
\item \emph{$q\le25$ is out of reach of this route.} At $q=25$, $M_{14}(25)\approx-0.0665<0$, and $14$ is the largest
budget available for this module (order $13$ plus the constant). Neither the rank nor the irreducibility can help; a new
family of forms, a better denominator type or a better map would be needed.
\item \emph{$10^{-5}$ is not reached.} For Theorem~\ref{thm:Cprime} the best range margin we found at $4\cdot10^{-6}$
is $-0.013$ (partition form), and even for a single pair the best margin at $10^{-5}$ is $-0.0086$ [numerical, local
optima]. We have no certified ceiling: the only rigorous ceiling we know is far from tight.
This agrees with \cite[Remark~14.0.4]{CDT}, which says that the constant could be improved but probably not beyond
about $10^{-4}$.
\item \emph{Pooling stops near $6\cdot10^{-5}$ for $K=2$} on the maps we explored (margins $-0.009$ at $6\cdot10^{-5}$,
$-0.031$ at $7\cdot10^{-5}$, $-0.069$ at $10^{-4}$, numerical).
\item \emph{The main inherited assumption}, the irreducibility of $\mathcal L_{a,b}$ (asserted in CDT's template
Lemma~12.1.1 and implicit in Lemma~14.3.1), is open. Its local system is the push-forward of a rank-$2$ system along the
double cover $x\mapsto y$; by Mackey's criterion it would follow from the irreducibility of the rank-2 system and its
non-isomorphism with its conjugate under $x\mapsto x/(x-1)$ [heuristic].
\end{itemize}

% ===========================================================================
\section{Provenance and use of AI tools}\label{sec:provenance}
% ===========================================================================

\subsection*{Provenance}
The mathematical inputs are: Lysov's explicit formula for $Q_n$ \cite{Lys18}, used as a definition (his Lemmas~7--9 are
re-proved above); the holonomy bounds and the construction for products of two logarithms of Calegari, Dimitrov and Tang
\cite{CDT}, used as stated; Hata's theorem \cite{Hat98}, used only for the extension to $q\ge54$ and for
Theorem~\ref{thm:F2}(iv). The scheme of Section~\ref{sec:proof4T} is the one suggested in \cite[\S15.1.10]{CDT}. The
recurrence of Proposition~\ref{prop:I1} was found by modular linear algebra and then proved by creative telescoping; the
minimal operators were found the same way. Apart from the preprints cited above (\cite{CDT}, \cite{Lys18} and, for
Lemma~\ref{lem:H}, \cite{LSS20}), nothing in this paper depends on unpublished work by others.

\subsection*{Use of AI tools}
This work was carried out with substantial assistance from AI systems (Anthropic's Claude models), which were used to
search the literature, to develop the proofs, to write the scripts and produce the certificates, to referee the
arguments, and to draft this text. The referee reports were also produced by Claude models, working adversarially on
separate copies of the files; they are not a substitute for a human expert's reading. The author directed the project
and is responsible for its content. None of the results is verified in Lean or another proof assistant.

% ===========================================================================
\bibliographystyle{amsplain}
\bibliography{refs}

\end{document}
